\documentclass[11pt]{article}
\usepackage[T1]{fontenc}
\usepackage{lmodern}
\usepackage[margin=26mm]{geometry}
\usepackage{amsmath,amssymb,amsthm,mathtools}
\usepackage{microtype,booktabs,longtable,array,xcolor}
\usepackage{enumitem,fancyhdr,listings}
\usepackage[colorlinks=true,linkcolor=blue!45!black,citecolor=blue!45!black,urlcolor=blue!45!black]{hyperref}
\usepackage{xurl}
\hypersetup{pdftitle={A substitution-stable seven-rule presentation of intuitionistic primitive recursive arithmetic},pdfauthor={Research prepared for Vladimir Reshetnikov},pdfsubject={Diagonal guard, exact native-kernel equivalence, and whole-domain bi-interpretation with iPRA},pdfkeywords={iPRA, primitive recursive arithmetic, intuitionistic logic, proof checker, bi-interpretation}}
\setlength{\headheight}{14pt}
\pagestyle{fancy}
\fancyhf{}
\fancyhead[L]{\small Concise iPRA: diagonal-guard introduction}
\fancyhead[R]{\small Research construction}
\fancyfoot[C]{\thepage}
\newtheorem{theorem}{Theorem}[section]
\newtheorem{lemma}[theorem]{Lemma}
\newtheorem{proposition}[theorem]{Proposition}
\newtheorem{corollary}[theorem]{Corollary}
\theoremstyle{definition}
\newtheorem{definition}[theorem]{Definition}
\theoremstyle{remark}
\newtheorem{remark}[theorem]{Remark}
\newcommand{\ipra}{\mathsf{iPRA}}
\newcommand{\FV}{\operatorname{FV}}
\newcommand{\one}{S0}
\lstset{basicstyle=\ttfamily\scriptsize,columns=fullflexible,breaklines=true,keepspaces=true,showstringspaces=false,frame=single,rulecolor=\color{black!20},xleftmargin=2mm,xrightmargin=2mm}
\setlist{topsep=4pt,itemsep=2pt}
\title{A substitution-stable seven-rule\texorpdfstring{\\}{ }presentation of intuitionistic\texorpdfstring{\\}{ }primitive recursive arithmetic}
\author{A research construction for Vladimir Reshetnikov}
\date{November 11, 2026}
\begin{document}
\maketitle
\begin{abstract}
The native scalar kernel in the current ConciseIPRA project has eight
primitive inference schemes. We replace its two introduction schemes by
one, retaining the complete original term and formula language and the
ordinary meanings of both binders. An atomic matrix receives its
progress hypothesis exactly when it has the syntactic form
\(I(y,y;u.h)=I(y,y;u.k)\); every other matrix receives its zero
instance. The zero instance of every recognized matrix is provable by
computation. This recovers all case introductions. A diagonal lift of an
arbitrary equation then recovers all equation induction. The resulting
seven-rule calculus derives exactly the original sequents and consequently
inherits the whole-domain bi-interpretation with full first-order iPRA.
The guard test is both preserved and reflected by capture-free substitution,
so the single rule supports the usual direct structural induction on proofs.

The matched, fully typed standalone Python checker has 3,244 normalized
bytes, 1,069 significant tokens and 1,124 AST nodes, compared with
3,382, 1,101 and 1,160 for the eight-rule control. This is a qualified
improvement: the existing witness-selected seven-rule control is smaller
at 3,205 bytes but its primitive mode is not stable under substitution.
The construction does not beat the 521-byte compact profiles listed in the pinned README
and does not establish absolute minimality. The new mathematics is proved
in writing; the accompanying Python certificates are checked, but no new
Lean or Rocq verification is claimed.
\end{abstract}
\tableofcontents
\newpage

\section{Result, comparison, and scope}
\label{sec:result-scope}

The comparison is frozen at ProveIt commit
\texttt{6bf83d8765614b0af51b35747e07fb057d6b2f42}, in
\texttt{Logic/PrimitiveRecursiveArithmetic/ConciseIPRA}~\cite{repo}.
Its README lists thirty-eight research packages, and its Lean development
proves the equivalence of the native scalar calculus, its ordinary
first-order realization, and iPRA~\cite{native,bridge}.

The result here is a new combination of properties: seven primitive rules,
the unrestricted original formula grammar, the unchanged conditional
readings of the binders, and an introduction rule whose mode is invariant
under capture-free numerical substitution. These are properties of the
axiomatization, not just of a Python encoding.

\begin{center}
\small
\begin{tabular}{@{}>{\raggedright\arraybackslash}p{34mm}c>{\raggedright\arraybackslash}p{34mm}>{\raggedright\arraybackslash}p{35mm}@{}}
\toprule
Presentation & Rules & Binder meanings & Introduction mode under substitution\\
\midrule
Native kernel & 8 & Original & Preserved; two rules\\
Witness-selected kernel & 7 & Original & Can change when a witness becomes literal zero\\
Uniform-progress kernel & 7 & Changed at equations & Preserved\\
Diagonal-guard kernel & 7 & Original & Preserved and reflected\\
\bottomrule
\end{tabular}
\end{center}

The prior witness rule selects progress from the matrix root and whether
the witness is literally zero~\cite{witness}. The prior uniform-progress
construction instead changes the reading of an atomic closure and modifies
closure elimination~\cite{uniform}. The prior guarded-core construction
restricts public atomic matrices beneath the paired binder to a guarded
shape~\cite{guarded}. The present construction requires no such public
restriction, changes no ordinary binder reading, and leaves closure
elimination untouched.

\paragraph{What ``simpler'' means in this report.}
There is a strict reduction from the repository's eight-rule formalized
kernel in rule count and in all three matched source-size coordinates.
There is also a structural advantage over its witness-selected seven-rule
alternative: substitution admissibility follows by a direct proof-tree
induction without changing introduction mode. There is no claim that this
candidate dominates every existing implementation, or that seven is a
lower bound. In particular, the 521-byte compact profiles listed in the
pinned README remain unbeaten. Those profiles assume a canonical guarded public goal and admit
administrative syntax, whereas the checkers measured here validate every
term, formula, scope, context and proof node in the unrestricted native
language. These are different input contracts.

\paragraph{Target theory.}
Throughout, iPRA means full first-order intuitionistic arithmetic with a
symbol for every fixed finite primitive-recursive definition, its defining
equations, the successor separation and injectivity laws, and
quantifier-free induction, with parameters and retained assumptions.
The atomic-induction formulation in the usual PR language is equivalent
through numerical characteristics of quantifier-free formulas. This is
not merely the quantifier-free equational fragment. The conventional
first-order usage is also discussed by Salehi~\cite{salehi}.

\section{A substitution-stable single introduction}
\label{sec:diagonal-core}

The construction retains every numerical term and every formula of the
native scalar calculus.  It also retains the ordinary readings of its two
binders.  The change is confined to the introduction rule: a particular
\emph{syntactic diagonal shape} selects progress; every other matrix selects
cases.  The diagonal is chosen so that its zero instance is a theorem and
its recognition commutes with capture-free substitution.

\subsection{Syntax, readings, and all seven primitive rules}

Numerical terms and native formulas are
\[
 t ::= x\mid 0\mid S(t)\mid I(n,a;u.h),
 \qquad
 A ::= (s=t)\mid Q_y A\mid D_{x,y}A.
\]
In an iterator only $u$ is bound, and only in $h$; in $D$ the variables
$x,y$ are distinct.  We work up to renaming of bound variables, and all
substitutions avoid capture.  No restriction is imposed on public formulas,
intermediate formulas, iterator bodies, or witnesses.  The designated
contradiction remains $0=S0$.  The ordinary readings remain
\[
 Q_y B := B(0)\to\forall y\,B(y),\qquad
 D_{x,y}B := \exists x\bigl(B(x,0)\to\forall y\,B(x,y)\bigr).
\]

\begin{definition}[The diagonal predicate]
For a matrix $B(x,y)$, let $\operatorname{Diag}_y(B)$ mean that its
\emph{literal syntax} is
\begin{equation}
 B(x,y)=
 \bigl(I(y,y;u.h)=I(y,y;v.k)\bigr)
 \label{eq:diag-shape}
\end{equation}
for some bodies $h,k$.  The displayed counter and initializer on both sides
are each the variable bound as $y$ by the impending $D$.
The state binders $u,v$ may be renamed independently.  The bodies may use
$x,y$, their own state variable, and any fixed parameters.
This is syntactic recognition, without simplification or an equality test
inside the arithmetic theory.  Put
\[
 \eta_B(z)=
 \begin{cases}
 z,&\operatorname{Diag}_y(B),\\
 0,&\text{otherwise}.
 \end{cases}
\]
\end{definition}

The calculus $\mathcal G$ has the following complete primitive inventory.
Contexts are finite reusable lists of native formulas.
\begin{align*}
 \frac{A\in\Gamma}{\Gamma\vdash A}
   &\quad\text{Hyp},\\[3pt]
 \Gamma\vdash I(0,a;u.h)=a
   &\quad\text{I0},\\[3pt]
 \Gamma\vdash I(Sn,a;u.h)=h[I(n,a;u.h)/u]
   &\quad\text{IS},\\[3pt]
 \frac{\Gamma\vdash s=t\qquad
       \Gamma\vdash a(s)=b(s)}
      {\Gamma\vdash a(t)=b(t)}
   &\quad\text{Tr}_{=},\\[6pt]
 \frac{\Gamma\vdash Q_y B(y)\qquad\Gamma\vdash B(0)}
      {\Gamma\vdash B(t)}
   &\quad\text{QE},\\[6pt]
 \frac{\Gamma\vdash D_{x,y}B(x,y)\qquad
       \Gamma,Q_yB(b,y)\vdash A}
      {\Gamma\vdash A}
   &\quad\text{DE},\\[6pt]
 \frac{\Gamma,B(t,\eta_B(z))\vdash B(t,Sz)}
      {\Gamma\vdash D_{x,y}B(x,y)}
   &\quad\text{DI}_{\mathcal G}.
\end{align*}
In $\mathrm{Tr}_{=}$, $a(v),b(v)$ are numerical term templates; transport
inside a compound formula is not primitive.  In DE, $b$ is fresh for
$\Gamma,A$ and the fixed parameters of $B$.  In DI, $z$ is fresh for
$\Gamma,t$ and the fixed parameters of $B$.  The displayed local
assumptions are discharged.  There is no primitive reflexivity,
substitution, cut, ordinary connective, case, or additional induction
rule.

For comparison, $\mathcal N$ is the eight-rule native kernel of the
repository: its first six rules are the ones displayed above, and its two
introductions are
\begin{equation}
 \frac{\Gamma,B(t,0)\vdash B(t,Sz)}
      {\Gamma\vdash D_{x,y}B(x,y)}
 \quad\text{DI}_{\rm case},
 \qquad
 \frac{\Gamma,E(z)\vdash E(Sz)}
      {\Gamma\vdash Q_y E(y)}
 \quad\text{QI}_{\rm eq},
 \label{eq:diag-native-intros}
\end{equation}
where $E$ is a literal equation and the same freshness conventions apply.

\subsection{Why the diagonal survives substitution}

\begin{lemma}[Preservation and reflection of the guard]
\label{lem:diag-stable}
Let $\sigma$ be a capture-free substitution for the free parameters of
$B(x,y)$, with $x,y$ retained as the distinguished bound coordinates.
Then
\[
 \operatorname{Diag}_y(B\sigma)
 \quad\Longleftrightarrow\quad
 \operatorname{Diag}_y(B).
\]
The same assertion holds when a witness is substituted for $x$, provided
that the witness does not contain the bound coordinate $y$.
\end{lemma}
\begin{proof}
Rename binders away from every term in the substitution range.
A constructor-headed term remains headed by that constructor after
substitution.  Moreover, a term becomes the distinguished bound variable
$y$ if and only if it was already $y$: another retained variable stays
different from $y$, a substituted parameter has an image not containing
$y$, and a constructor-headed term cannot become a variable.
Consequently an iterator becomes headed by the exact pair of arguments
$(y,y)$ if and only if it had that pair before substitution.
The root of a formula is also preserved.  An entire side that is a
substituted parameter may become iterator-headed, but cannot acquire the
specified shape, because its image has no free occurrence of the protected
$y$.  An existing iterator matches the shape exactly when its two displayed
arguments match $y$, by the preceding argument.  Applying these facts to
both sides proves both directions of the equivalence.  The bodies are
unrestricted, so they impose no further condition.  Witness substitution
is the same argument with $x$ among the substituted parameters.
\end{proof}

The reflection direction is essential.  The superficially similar guard
\[
 I(y,0;u.h)=I(y,0;v.k)
\]
does not have it.  For a free parameter $p$, the matrix
\[
 I(y,p;u.0)=I(y,0;v.0)
\]
does not match the zero-initializer guard, whereas its substitution
$p:=0$ does.  Thus that test can change a case instruction into a progress
instruction.  Requiring the initializer itself to be the protected bound
variable $y$ prevents this event.

\begin{lemma}[Structural admissibility]
\label{lem:diag-structural}
Weakening, exchange, contraction, renaming, capture-free numerical
substitution, and ordinary-assumption cut are admissible in $\mathcal G$.
Numerical substitution can be performed by induction on a derivation
without changing any last-rule constructor or either mode of
$\mathrm{DI}_{\mathcal G}$.
\end{lemma}
\begin{proof}
Weakening, exchange and contraction change only hypothesis membership;
rename eigenparameters before inserting a context that could contain
them.  For numerical substitution, induct on the derivation.  The
computation and transport rules commute with capture-free substitution.
For QE and DE use the usual interchange of lifted substitution and
instantiation, after renaming their bound and eigenvariables.  In DI,
Lemma~\ref{lem:diag-stable} says that the same one of $0,z$ is retained
after substitution.  The substituted premise is therefore exactly the
premise of the same mode of DI.  This proves the induction step directly.

For cut, replace each use of a discharged ordinary hypothesis by its
supplied derivation.  Induct on the derivation in which that hypothesis
occurs, weakening the inserted proof in each local context.  Rename local
eigenparameters away from the inserted proof.  The guard is independent
of context and proof terms, so its mode is unchanged.  This proves both
single-assumption cut and simultaneous replacement of hypotheses by
derivations.  None of these arguments appeals to the equivalence with
$\mathcal N$ proved below.
\end{proof}

\subsection{Recovering every case introduction first}

\begin{lemma}[Elementary equality and the diagonal base]
\label{lem:diag-base}
Reflexivity, symmetry, transitivity and congruence for numerical terms are
derivable using only I0 and $\mathrm{Tr}_{=}$.  If
$\operatorname{Diag}_y(B)$, then $\Gamma\vdash B(t,0)$ for every
context $\Gamma$ and witness $t$.
\end{lemma}
\begin{proof}
Put $r=I(0,a;u.0)$.  The computation $r=a$, used as both transport
premises with equation template $v=a$, gives $a=a$.
From $s=t$ and $s=s$, the template $v=s$ gives $t=s$.
The template $s=v$ gives transitivity, and the template
$h(s)=h(v)$ gives congruence after reflexivity.
At $y=0$, both sides of~\eqref{eq:diag-shape} are of the form
$I(0,0;u.h')$, hence each equals $0$ by I0.
Symmetry and transitivity give their equality.  No induction, body
extensionality, or test for numerical equality is used.
\end{proof}

\begin{lemma}[Unrestricted case introduction]
\label{lem:diag-all-cases}
Every instance of $\mathrm{DI}_{\rm case}$ in
\eqref{eq:diag-native-intros} is admissible in $\mathcal G$, for every
native matrix $B$.
\end{lemma}
\begin{proof}
If $B$ does not match the diagonal, the desired rule is primitive.
Otherwise start with
$\Gamma,B(t,0)\vdash B(t,Sz)$.
By Lemma~\ref{lem:diag-base} the discharged base is already a theorem,
so cut it out.  This gives $\Gamma\vdash B(t,Sz)$.
Weaken by $B(t,z)$ and apply the progress mode of
$\mathrm{DI}_{\mathcal G}$.
This recovers the omitted case instance without any general equation
induction.
\end{proof}

\begin{lemma}[Conditional closure and ordinary numerical cases]
\label{lem:diag-case-interface}
For every native formula template $F(y)$, the rules
\[
 \frac{\Gamma,F(0)\vdash F(Sz)}{\Gamma\vdash Q_yF(y)},
 \qquad
 \frac{\Gamma\vdash F(0)\qquad\Gamma,F(0)\vdash F(Sz)}
      {\Gamma\vdash F(t)}
\]
are admissible in $\mathcal G$, with $z$ fresh for the retained
context and fixed parameters.  In particular, the second successor
premise may ignore its additional hypothesis $F(0)$.
\end{lemma}
\begin{proof}
Choose an unused witness coordinate $x$.  The recovered case rule at
witness $0$ introduces $D_{x,y}F(y)$ from the first displayed premise.
DE opens $Q_yF(y)$, which is independent of the fresh witness; return
that hypothesis as the conclusion of DE.  This proves the first rule.
QE with the base proof proves the second rule.
\end{proof}

\subsection{A small logical tool with no hidden induction}

Only implication between two equations is needed in the completeness
proof.  It can be constructed directly, so no general formula selector,
reflected equality, predecessor, or pairing theorem is needed at this
stage.  Define the numerical selector
\[
 C(w,a,b):=I(w,a;u.b),\qquad u\notin\operatorname{FV}(b).
\]
Computation gives $C(0,a,b)=a$ and $C(Sw,a,b)=b$.
For equations $U=(a=b)$ and $V=(c=d)$, with $w$ fresh, put
\begin{equation}
 U\Rightarrow_{=}V
   :=Q_w\bigl(C(w,a,c)=C(w,b,d)\bigr).
 \label{eq:diag-equational-implication}
\end{equation}
This is a native formula abbreviation, not an additional connective.

\begin{lemma}[Implication for equations]
\label{lem:diag-eq-implication}
For literal equations $U,V$, the rules
\[
 \frac{\Gamma,U\vdash V}{\Gamma\vdash U\Rightarrow_{=}V},
 \qquad
 \frac{\Gamma\vdash U\Rightarrow_{=}V\qquad\Gamma\vdash U}
      {\Gamma\vdash V}
\]
are derivable using computation, equational transport, and the
case-derived closure rule of Lemma~\ref{lem:diag-case-interface}.
\end{lemma}
\begin{proof}
Let $J(w)$ denote the equation inside the $Q$ in
\eqref{eq:diag-equational-implication}.
Its zero instance is interderivable with $U$, and each successor
instance is interderivable with $V$, by the two selector computations
and elementary equality reasoning.
For introduction, convert the retained $J(0)$ to $U$, insert the given
derivation of $V$, convert $V$ to $J(Sz)$, and apply case-derived
$Q$ introduction.  For elimination, convert the proof of $U$ to
$J(0)$, apply QE at $S0$, and convert $J(S0)$ to $V$.
All conversions are between equations.  This argument uses neither
compound transport nor general progress.
\end{proof}

\begin{lemma}[Cases in a context that mentions the inspected term]
\label{lem:diag-dependent-cases}
Let $E$ be an equation and let $r$ be any term.  The rule
\[
 \frac{\Gamma,r=0\vdash E\qquad\Gamma,r=Sv\vdash E}
      {\Gamma\vdash E}
\]
is admissible when $v$ is fresh for $\Gamma,r,E$.
The context $\Gamma$ may contain $r$ and its free variables.
\end{lemma}
\begin{proof}
Keep $\Gamma,r,E$ fixed and choose a new variable $y$.
Set $F(y):=((r=y)\Rightarrow_{=}E)$.
Discharge the two branch equations by
Lemma~\ref{lem:diag-eq-implication}; this gives proofs of $F(0)$ and
$F(Sv)$ in $\Gamma$.  Apply the ordinary case rule of
Lemma~\ref{lem:diag-case-interface} to $F(y)$ at target $r$.
The result is $(r=r)\Rightarrow_{=}E$.
Reflexivity and implication elimination give $E$.
The variable being generalized in the case inference was the fresh
$y$, not a variable occurring in $\Gamma$.  Thus the proof does not
use the invalid convention of changing a context silently in each branch.
\end{proof}

\subsection{Recovering unrestricted equation progress}

\begin{theorem}[The full native progress rule]
\label{thm:diag-progress}
For every literal equation $E(y)$, the rule
\[
 \frac{\Gamma,E(z)\vdash E(Sz)}{\Gamma\vdash Q_yE(y)}
\]
is admissible in $\mathcal G$, with the usual eigenparameter condition.
\end{theorem}
\begin{proof}
Write $E(y)=(s(y)=t(y))$ and let $\pi$ be the supplied step derivation.
Define the diagonal lift
\begin{equation}
 H(y):=
 \bigl(I(y,y;u.s(y))=I(y,y;u.t(y))\bigr),
 \qquad u\notin\operatorname{FV}(s,t).
 \label{eq:diag-lift}
\end{equation}
By computation and equality alone,
\begin{equation}
 \vdash H(0),\qquad
 H(Sv)\dashv\vdash E(Sv)
 \label{eq:diag-lift-endpoints}
\end{equation}
for every term $v$.  The second statement uses constant state updates;
it does not assert any identity between arbitrary iterator bodies.

Work temporarily in the retained context
$\Delta:=\Gamma,E(0)$, which contains no fresh induction variable.
For fresh $y$ set
\[
 \Phi(y):=H(y)\Rightarrow_{=}E(Sy).
\]
There are two case proofs in this fixed context:
\begin{itemize}
\item At $y=0$, substitute $0$ for the eigenparameter of $\pi$.
  The retained $E(0)$ gives $E(S0)$.
  The assumption $H(0)$ may be ignored, so equation-implication
  introduction gives $\Delta\vdash\Phi(0)$.
\item At $y=Sv$, assume $H(Sv)$.
  By~\eqref{eq:diag-lift-endpoints} it gives $E(Sv)$.
  Substitute $Sv$ for the eigenparameter of $\pi$ and obtain $E(SSv)$.
  Discharging $H(Sv)$ gives $\Delta\vdash\Phi(Sv)$.
\end{itemize}
These uses of substitution are justified directly by
Lemma~\ref{lem:diag-structural}; they do not invoke the result currently
being proved.
The ordinary case rule applies to the \emph{compound} native formula
$\Phi$ and yields $\Delta\vdash\Phi(z)$ for a fresh $z$.
Now assume $H(z)$, use equation-implication elimination, and convert
$E(Sz)$ back to $H(Sz)$.  We have proved
\[
 \Delta,H(z)\vdash H(Sz).
\]
The case inference was performed \emph{before} adding $H(z)$ to the
context, so its freshness condition is satisfied.  Equivalently, one
could use the explicitly justified contextual case rule of
Lemma~\ref{lem:diag-dependent-cases}.

Take an unused witness coordinate $x$ in the matrix $H(y)$.
It has the exact diagonal shape, so primitive progress DI at witness
$0$ gives $\Delta\vdash D_{x,y}H(y)$.
DE, with its fresh witness unused, gives $\Delta\vdash Q_yH(y)$.
Together with $H(0)$, QE at a fresh successor $Sz$ yields $H(Sz)$,
hence $E(Sz)$.  We have therefore constructed
\[
 \Gamma,E(0)\vdash E(Sz).
\]
Finally apply the already recovered case-derived $Q$ introduction to
obtain $\Gamma\vdash Q_yE(y)$.

The temporary retention and final discharge of $E(0)$ are deliberate.
No inference identifies conditional closure with unconditional progress.
In particular, the argument also works when the base equation is false.
\end{proof}

\subsection{Exact equivalence and the arithmetic consequence}

\begin{theorem}[Identity-language equivalence]
\label{thm:diag-equivalence}
For every finite context $\Gamma$ and every native formula $A$,
\[
 \boxed{\quad
 \Gamma\vdash_{\mathcal G}A
 \quad\Longleftrightarrow\quad
 \Gamma\vdash_{\mathcal N}A.
 \quad}
\]
Both translations act effectively on finite derivations, and neither
translation changes a term, formula, or retained assumption.
\end{theorem}
\begin{proof}
Translate a $\mathcal G$ derivation to $\mathcal N$ by induction.
The first six rules are shared.  An unguarded DI instance is native
case introduction.  In a guarded DI instance the matrix is an equation.
The native progress rule applied to $E(y):=B(t,y)$ gives $Q_yB(t,y)$.
Under the temporary base $B(t,0)$, QE supplies $B(t,Sz)$; native case
DI then introduces the required $D_{x,y}B(x,y)$.

Conversely, translate a native derivation recursively.
The shared rules remain unchanged, every native case introduction is
replaced by Lemma~\ref{lem:diag-all-cases}, and every native equation
progress introduction is replaced by
Theorem~\ref{thm:diag-progress}.  All constructions use the original
retained context and satisfy the freshness conditions.
\end{proof}

The proof has a strict dependency order:
guard stability; structural admissibility; elementary equality;
all case introductions; case-derived closure and numerical cases;
implication between equations; general equation progress.
It does not use equality reflection, a scalar predecessor, pairing,
primitive-recursion compilation, or the eventual bi-interpretation in
order to establish equivalence.

\begin{proposition}[Transfer of the established bi-interpretation]
\label{prop:diag-ipra}
Compose Theorem~\ref{thm:diag-equivalence} with the native scalar
calculus's established bi-interpretation with iPRA.  This gives a
bi-interpretation of $\mathcal G$ with iPRA, and equiconsistency.
The numerical domains and interpreted equality are unchanged, and both
numerical composites have the identity graph $x=y$.
\end{proposition}
\begin{proof}
The first component of each composite, between $\mathcal G$ and
$\mathcal N$, is literally the identity on the full syntax and on all
sequents.  It therefore preserves the established whole-domain
interpretations, their internal formula round trips, and their numerical
identity graphs.  The equivalence also preserves the designated native
contradiction $0=S0$ in the empty context.  Composing the finite proof
translations in each direction preserves, and reflects, derivability of
contradiction, which proves equiconsistency.
\end{proof}

\subsection{A sufficient criterion for other guards}

\begin{proposition}[The guard criterion]
\label{prop:diag-guard-criterion}
Let $\mathsf P_y(B)$ be a decidable predicate on native matrices that
respects capture-free renaming.  Suppose that:
\begin{enumerate}
\item Every matrix satisfying $\mathsf P$ is a literal equation, and
  its instance $B(t,0)$ is derivable for every witness $t$ using only
  computation and equational transport.
\item Every equation $E(y)$ has an equation lift $H(y)$, independent
  of an added unused witness coordinate $x$, such that
  $\mathsf P_y(H)$ holds and
  $H(Sv)\dashv\vdash E(Sv)$ is derivable using only computation
  and equational transport, for arbitrary $v$.
\end{enumerate}
Replace the diagonal test in the seven-rule presentation by $\mathsf P$.
The resulting calculus $\mathcal G_{\mathsf P}$ derives exactly the
native sequents.  If, in addition, $\mathsf P$ is preserved and reflected
by capture-free numerical substitution, its structural substitution
theorem follows by direct rule-preserving induction.
\end{proposition}
\begin{proof}
The first condition recovers every missing case introduction exactly as
in Lemma~\ref{lem:diag-all-cases}.  Hence the case-derived closure,
numerical cases and equation-implication rules are available.
The second condition supplies the only properties of the lift used in
Theorem~\ref{thm:diag-progress}.  Since the base and endpoint conversions
belong to the fragment with computation and equational transport, neither
condition presupposes the general progress rule being recovered.

For clarity, equivalence does not silently assume numerical substitution
in $\mathcal G_{\mathsf P}$ when the optional stability condition is
absent.  Translate native derivations by strong induction on their height.
In the native progress case, substitute $0$ and $Sv$ into its native
premise derivation first.  Native substitution preserves derivation
height, so both resulting proofs still have smaller height than the
progress conclusion and can be translated by the induction hypothesis.
They supply the two $\Phi$ branches in the displayed proof above.
Renaming, weakening and cut require only the stated renaming invariance.
The other native rules translate as before, and the reverse translation
uses only the fact that a progress-selected matrix is an equation.
This proves the sequent equivalence.  Under the additional preservation
and reflection condition, numerical substitution instead follows directly
from the proof of Lemma~\ref{lem:diag-structural}, with no change of mode.
\end{proof}

\subsection{What is simpler, and what is not being claimed}

Relative to the eight-rule native kernel, this construction removes one
primitive introduction while retaining its full language, all witnesses,
and both original binder readings.  Relative to the earlier seven-rule
witness-selected presentation, its improvement is structural: the mode
depends only on protected bound-variable positions in the matrix, so
ordinary substitution is proved by a direct rule-preserving induction.
There is no exceptional substitution-repair proof.

The earlier guarded-language proposal restricted the public atomic
matrices; here every original formula remains available and the comparison
theorem is the identity on \emph{all} sequents.  The earlier
uniform-progress proposal changes the meaning of $Q$ at equations and
modifies its elimination; here neither change is made.  The diagonal
initializer supplies the missing combination of a provable zero base and
a substitution-invariant test.

This is not a proof of global minimality.  A seven-rule count is not a
new record by itself, and the syntactic guard has an implementation cost.
Any claim about a smaller complete checker must be supported by the
separate matched source measurements, including syntax validation,
substitution, freshness, and the public certificate interface.
The new equivalence above is a written mathematical proof; it is not, by
itself, a claim that a Lean or Rocq verification of this new kernel has
been compiled.


\section{Ordinary logic and the bi-interpretation with iPRA}
\label{sec:interpretation}

The equality of complete consequence relations established above is what
allows the existing arithmetic bridge to be reused. Agreement on closed
numerals or on equational theorems alone would not suffice.

\subsection{Ordinary realization and explicit logical content}

Let \(T_I\) be intuitionistic first-order logic with actual equality, the
scalar terms of the new calculus, its two iterator computation schemes,
and the following additional arithmetic principles:
\[
 \forall x\,(x=0\lor\exists z\,x=Sz),\qquad
 (0=S0)\to\bot,
\]
together with contextual induction for literal equations. These are
principles of the ordinary realization; in the native seven-rule calculus
they are derived, as are the usual logical inference rules.
The realization is
\[
 (s=t)^\circ=(s=t),\quad
 (Q_y B)^\circ=B^\circ(0)\to\forall y\,B^\circ(y),
\]
\[
 (D_{x,y}B)^\circ=
 \exists x\bigl(B^\circ(x,0)\to\forall y\,B^\circ(x,y)\bigr).
\]
All interpretation domains are the whole numerical domain, defined by
\(x=x\). No quotient equality is used.

For clarity, there are no unstated primitive logical rules in the new
calculus. Their native definitions may be written using the selector
\(C(v,a,b)=I(v,a;u.b)\), a formula selector \(M_v(A,B)\), and a formula
gate \(G_vA\). At equations the selector applies \(C\) to corresponding
sides and the gate has zero on both sides at its zero branch. At binders
these operations recurse through the formula; unmatched binder depths are
padded using an equivalent dummy-witness gated formula. First replace any
primitive \(Q_yB\) by the interderivable \(D_{x,y}B\) with an unused
witness coordinate \(x\); the case-derived closure rules prove this in
both directions. The explicit recursive definitions on the resulting
Eq/D syntax and their endpoint proofs are reproduced in
Appendix~\ref{sec:choice-introduction}.

With \(p(v)=I(v,0;u.C(u,1,0))\), one may use the following full set of
macros (fresh bound variables are always chosen):
\begin{align*}
 \bot_*&=(0=S0),\\
 A\Rightarrow_* B&=D_{u,z}M_z(A,B),\\
 A\wedge_* B&=D_{u,z}G_zM_{p(z)}(A,B),\\
 \exists_*x\,A(x)&=D_{x,y}G_yA(x),\\
 A\vee_*B&=\exists_*x\,M_x(A,B),\\
 \forall_*x\,A(x)&=A(0)\wedge_*D_{u,z}G_zA(z).
\end{align*}
The variable \(u\) is unused. The case rules give all ordinary
intuitionistic introduction and elimination rules for these macros.
The progress rule supplies exactly the additional equation induction.
No excluded-middle rule for arbitrary formulas is added.

The required separation sentence in \(T_I\) must not be omitted:
positive equations alone admit a singleton model. In the native reading,
\(0=S0\) is the designated contradiction and explosion from it is
derivable by the selector equations. The ordinary separation sentence
is precisely what connects that designated equation with ordinary
first-order falsity.

\subsection{The two term interpretations}

An iterator in iPRA is defined by primitive recursion with an update that
ignores the recursion counter:
\[
 i_h(0,a,\vec p)=a,\qquad
 i_h(Sn,a,\vec p)=h(i_h(n,a,\vec p),\vec p).
\]
This is an external translation on each fixed finite term description;
there is no internal universal evaluator for all primitive-recursive
programs.

In the other direction the scalar theory first constructs predecessor,
then a faithful equality characteristic, then quantifier-free induction,
and then pairing and projections. The dependency order matters: pairing
and equality reflection cannot be used as assumptions in the predecessor
construction. Complete inherited arguments are reproduced in
Appendices~\ref{sec:hybrid-bootstrap}, \ref{sec:eic-completion} and
\ref{sec:iterator-coding}.

If \(\pi,L,H\) are the resulting pairing and projections, a fixed PR
definition
\[
 f(0,\vec p)=g(\vec p),\qquad
 f(Sn,\vec p)=h(n,f(n,\vec p),\vec p)
\]
is compiled as
\begin{align*}
 V(n,\vec p)&=I\bigl(n,\pi(0,g^*(\vec p));
 u.\pi(SL(u),h^*(L(u),H(u),\vec p))\bigr),\\
 f^*(n,\vec p)&=H(V(n,\vec p)).
\end{align*}
Equation induction proves \(L(V(n,\vec p))=n\), and the projection
identities prove the defining equations of \(f^*\).

\subsection{Both composites and consistency}

For each fixed iterator and each fixed PR definition, equation induction
proves equality of two fixed term realizations satisfying the same
defining equations.
External induction on finite descriptions therefore identifies every
operation with its double translation. The numerical composite maps are
the identity and are internally represented by \(x=y\).

For formulas, the original native-to-ordinary and ordinary-to-native
translations already have both provable round trips. The substantive
paired-binder round trip is
\[
 D_{x,y}B\quad\dashv\vdash\quad
 \exists_*x\bigl(B(x,0)\Rightarrow_*\forall_*y\,B(x,y)\bigr).
\]
The closure round trip is obtained using a dummy witness. Because the new
calculus has exactly the original sequents, every such round-trip proof
is also a proof in the new calculus. Thus the composition with the
\(T_I\)--iPRA bridge is a genuine bi-interpretation on the whole
numerical domain, not merely a mutual consistency reduction.

Effective transformations of finite derivations carry a proof of native
\(0=S0\) to an iPRA proof of \(\bot\), and conversely. This yields
equiconsistency. It asserts neither that iPRA proves its own consistency
nor that the new system has a special standard-domain semantics.

\section{Standalone checkers and reproducible measurements}
\label{sec:implementation}

The three compared checkers share the same explicit de Bruijn syntax,
ordinary JSON input, scope checks, equality rule and input/output contract.
They differ only in their introduction rules. The eight-rule control is
the repository's \texttt{one\_context/direct/native\_checker.py} at the
pinned commit; the witness and diagonal variants are matched edits of that
control. The filename of the control records its package location, not a
change to the two-template equality rule used in this comparison.

The data object has exactly four fields:
\begin{lstlisting}
{"variables": n, "context": [...], "goal": formula, "proof": proof_tree}
\end{lstlisting}
Here \(n\) is a nonnegative integer. Context entries and the goal must be
well-scoped formulas. Terms are tagged as \texttt{v}, \texttt{z},
\texttt{s}, and \texttt{i}; formulas as \texttt{=}, \texttt{q} and
\texttt{d}. Only an iterator's body binds its state variable.
The matrix of \(D\) binds \(y\) at index 0 and \(x\) at index 1.
All proof nodes are checked and synthesize their conclusions.

The new introduction is implemented by opening the witness and retaining
the opened equation itself exactly when both iterator prefixes are
\texttt{['i', ['v', 0], ['v', 0]]}; otherwise it retains the shifted
zero instance. Well-formedness is checked before testing these prefixes.
Witness substitution cannot manufacture the bound \(y\) in one of these
positions, so checking the guard after opening the witness is equivalent
to checking the original matrix.

Acceptance means exit code zero and exactly \texttt{ACCEPT} followed by a
newline. The programs use neither assertions for logical checks nor
\texttt{eval}, a proof database, an external prover, a trusted compiler,
or an imported axiom file. The untrusted proof compiler and test producer
are separate from the checker. The trusted runtime includes CPython and
its standard \texttt{json} and \texttt{sys} modules, just as the controls
do. Resource exhaustion may reject a large valid finite certificate;
no claim of fixed-machine resource completeness is made.

\paragraph{Rule correspondence.}
The implementation uses one branch for each displayed primitive rule.
By induction on a successful call to \texttt{infer}, its result is a
well-scoped formula derivable from the supplied context: the hypothesis
branch checks membership, the computation branches check their term
arguments, and transport and QE compare the exact substituted minor
premise. DE checks its continuation against a shifted conclusion under a
new variable, excluding eigenvariable escape. DI checks exactly the
single local hypothesis selected by the guard. The binder increments in
\texttt{walk} are zero for ordinary term arguments, one for an iterator
body and a Q matrix, and two for a D matrix; structural induction gives
the capture-free meanings of \texttt{shift} and \texttt{sub}.
Conversely every finite derivation of the displayed calculus has the
corresponding tagged certificate. This is a written correspondence
argument for the idealized unbounded-resource program, not a
proof-assistant-checked theorem about the Python implementation.

\begin{center}\small
\begin{tabular}{@{}lrrrr@{}}
\toprule
Complete checker & Bytes & Tokens & AST nodes & gzip-9\\
\midrule
Native, eight rules & 3,382 & 1,101 & 1,160 & 990 \\
Witness-selected, seven rules & 3,205 & 1,041 & 1,099 & 975 \\
Diagonal-guard, seven rules & 3,244 & 1,069 & 1,124 & 989 \\
\bottomrule
\end{tabular}\end{center}

The normalization applies \texttt{ast.unparse(ast.parse(source))} and
appends one line-feed byte, under CPython 3.12.14. Significant tokens exclude encoding, indentation,
dedentation, newlines, comments and the end marker. AST nodes are counted
with \texttt{ast.walk}. Informational gzip sizes use compression level 9
and \texttt{mtime=0}. The same procedure is used for all three modules;
none of these figures omits a checker helper.

The diagonal rule saves 138 normalized bytes (4.08\%), 32 tokens
(2.91\%) and 36 AST nodes (3.10\%) against the eight-rule control.
It costs 39 bytes against the witness-selected control. This is a
Pareto tradeoff between source size and a substitution-invariant primitive
interface, not a new record for unrestricted minimum source size.

\section{Verification and limitations}
\label{sec:validation}

The reproducible implementation checks used 15 named proof fixtures and 17 additional malformed or scope-attack fixtures. Every fixture ran against all three standalone checkers, both normally and under \texttt{python -O}: 192 processes in total, with zero mismatches. The positive fixtures include a diagonal-guard induction whose hypothesis is used, the derived identity closure $Q_y(I(y,0;u.Su)=y)$, nonvacuous existential elimination, compound $Q$ and $D$ matrices, and a nested capture-avoiding witness. Expected rejections distinguish ordinary atomic cases from progress, reject a premature eigenvariable, and reject unrestricted compound induction.

There were also 240 finite checks of guard preservation and reflection under parameter substitution or witness opening, together with an explicit counterexample to stability of the older zero-seed guard. Independent standard-natural-number evaluation confirmed 80 capture-free term substitutions, with no mismatches or budget skips. These tests check the implementation and boundary conditions; they do not replace the written metatheory.


\paragraph{Executable translations.}
The untrusted translator implements both complete proof-tree translations,
including elimination of arbitrary native \texttt{qi} nodes. A separate
suite checks 29 certificate cases, with 93 checker validations covering
every source rule in both calculi and both round trips. In particular it
compiles an identity induction, an induction whose base is the unproved
equation \(a=b\), and an induction using both the current induction
hypothesis and a retained parameter equation. The three translated
proofs are also accepted through the public command-line interface in
ordinary and optimized Python modes. The count separates 87 direct calls
to the standalone checker modules from six command-line processes.

\begin{center}\small
\begin{tabular}{@{}lrr@{}}
\toprule
Certificate & Original JSON bytes & Translated JSON bytes\\
\midrule
Identity induction & 485 & 11,506\\
Conditional seed equality & 885 & 14,218\\
Contextual update congruence & 1,965 & 20,283\\
\bottomrule
\end{tabular}
\end{center}
These are certificate sizes, not checker sizes. The larger proof trees
make explicit the logical work removed from the primitive kernel. The
report claims an effective translation and reports these examples; it
does not claim a polynomial bound on serialized tree size.

The checker tests establish that the supplied finite certificates are
processed as reported. They do not establish universal soundness,
completeness, or consistency. Those claims rest on the written
metatheorems and the inherited arithmetic bridge. The pinned repository's
Lean formalization verifies the original bridge; it does not automatically
verify a new Python checker or the new diagonal-rule comparison.
No new Lean or Rocq build was performed for this delivery.

The archive includes the new construction, checker sources, proof
translation code and examples, measurements, reproduction scripts,
provenance, and the inherited arithmetic arguments with attribution.
The report's new contribution is the diagonal dispatch theorem and its
matched implementation. It does not claim a new predecessor construction,
a new PR compilation method, an absolute minimum number of rules, or a
smaller version of the 521-byte guarded-profile checker listed in the pinned README.

\section{Further research questions}
\label{sec:research}

\begin{enumerate}
\item \textbf{Lean comparison with the existing native kernel.}
Add a deeply embedded \(\mathcal G\) with the guard as a structural
predicate. Prove guard preservation and reflection first, followed by
substitution and cut directly. Then implement the equation-only
implication construction and the proof of unrestricted progress.
The target is an exact sequent theorem, not only consistency.
\item \textbf{The smallest stable guard.}
Can a cheaper decidable structural class retain all three required
properties: provable zero instances, a lift of arbitrary equations,
and reflection under substitution? The diagonal guard is a concrete
solution, not a lower-bound theorem.
\item \textbf{Proof size and sharing.}
The translations are effective, but syntactic lifting and repeated local
derivations may expand tree certificates. Determine explicit size bounds
and compare trees with a checked DAG format. Sharing must not enlarge the
trusted logical rule set implicitly.
\item \textbf{Combining independent equality reductions.}
The repository's crossed one-template transport is compatible at the
mathematical level with the present introduction theorem. A matched
implementation could determine whether their combination improves the
smallest fully typed checker while retaining a transparent proof format.
Such a measurement is separate from the results reported here.
\item \textbf{A fair compact-profile comparison.}
Can the new substitution-stable rule be encoded within a compact
administrative representation with a complete public-goal validator,
and improve its normalized source size? The current typed experiment
does not answer that question.
\end{enumerate}

\appendix
\section{Attribution and inherited proof material}
\label{sec:attribution}

The following five sections reproduce the mathematical bridge from the
pinned repository's \texttt{Research/certified\_zero/reference/} material
with light introductory wording changes. The Choice-DI theorem and its
logical macros, the ordinary realization, scalar predecessor bootstrap,
equality reflection, quantifier-free induction, pairing and PR compilation
are inherited results, not new results of this report. Exact fetched
sources and their provenance are included separately in the archive.
The repository distributes these materials under the MIT No Attribution
license; its license text is included.

% New mathematical result: complementary introduction choices.
% Requires amsmath/amssymb/amsthm and theorem, lemma, proposition environments.
% The inherited arithmetic-to-iPRA bridge must be supplied separately.

\section{Complementary introductions with one local assumption}
\label{sec:choice-introduction}

This inherited comparison starts from a paired introduction retaining both the zero instance and an
induction instance when its matrix is an equation. Both together are
unnecessary. A certificate may choose either one, provided the induction
choice is restricted to a literal equation. This changes neither the
formula language nor the meaning of a formula. The omitted combination
is recovered by a short nested introduction--elimination construction.

\subsection{Complete native syntax and rule inventory}

The scalar numerical language is
\[
 t::=x\mid0\mid St\mid I(n,a;u.h).
\]
Only $h$ binds the state variable $u$; $n$ and $a$ do not.
The update may contain fixed numerical parameters. The computation
axioms in every retained context are
\[
 I(0,a;u.h)=a,\qquad
 I(Sn,a;u.h)=h[I(n,a;u.h)/u].
 \tag{I0, IS}\label{eq:choice-computations}
\]
The argument also applies unchanged to the full fixed-description PR
language with its projection, composition, and recursion equations.
No numerical variable ranges over function descriptions, and no uniform
evaluator is postulated.

Native formulas and reusable context clauses are
\[
 A::=(s=t)\mid\mathsf D_{x,y}A,\qquad
 \langle A;y.B(y)\rangle,\quad y\notin\operatorname{FV}(A).
\]
Bound variables are renamed before capture-free substitution.
The two variables $x,y$ in a paired formula are distinct.
An ordinary assumption $F$ is the constant clause
$\langle H;v.F\rangle$, where $v$ is absent from $F$ and $H$ is the
fixed computation theorem $I(0,0;u.0)=0$. Thus ordinary hypothesis
use is a clause Use preceded by this computation proof.

There are computation axioms and exactly the following inference
patterns:
\[
\frac{\Gamma\vdash A}{\Gamma\vdash B(t)}
\quad\mathsf{Use}
\quad(\langle A;y.B(y)\rangle\in\Gamma),
\]
\[
\frac{\Gamma\vdash\mathsf D_{x,y}B(x,y)\qquad
      \Gamma,\langle B(a,0);y.B(a,y)\rangle\vdash F}
     {\Gamma\vdash F}
\quad\mathsf{DE},
\qquad
\frac{\Gamma\vdash s=t\qquad\Gamma\vdash F(s)}
     {\Gamma\vdash F(t)}
\quad\mathsf{Tr},
\]
\[
\frac{\Gamma,B(t,\eta_\varepsilon(z))\vdash B(t,Sz)}
     {\Gamma\vdash\mathsf D_{x,y}B(x,y)}
\quad\mathsf{DI}_{\varepsilon},
\qquad
\eta_0(z)=0,\quad\eta_1(z)=z.
\tag{Choice DI}\label{eq:choice-di}
\]
Here $\varepsilon$ is a certificate choice in $\{0,1\}$, and
$\varepsilon=1$ is permitted only when $B$ has a literal equation
root. An equation beneath a binder is not a literal equation.
Call $\varepsilon=0$ the \emph{case choice} and $\varepsilon=1$
the \emph{progress choice}. Both retain exactly one ordinary
formula assumption. Compound matrices have only the case choice.

In DI the eigenparameter $z$ is fresh for $\Gamma$, the witness
$t$, and every fixed parameter of $B$. In DE the eigenparameter
$a$ is fresh for $\Gamma$, the conclusion $F$, and every fixed
parameter of $B$. In transport, $F(v)$ is any fixed native
formula template. Displayed local assumptions and clauses are
discharged. There are no primitive reflexivity, coverage,
separation, ordinary logical, or additional induction rules.

We call this calculus $\mathcal C$. The original comparison
calculus $\mathcal N$ has identical syntax, computation,
Use, DE, and transport, but its introduction is
\[
\frac{\Gamma,B(t,0),[B(t,z)]_{\mathrm{at}}\vdash B(t,Sz)}
     {\Gamma\vdash\mathsf D_{x,y}B(x,y)}.
\tag{Old DI}\label{eq:choice-old-di}
\]
The bracket supplies $B(t,z)$ exactly at a literal equation.

\subsection{Exact equivalence and the unused-witness construction}

\begin{lemma}[Structural admissibility]
\label{lem:choice-structural}
Weakening, exchange, contraction, renaming, ordinary-assumption cut,
and capture-free numerical substitution are admissible in $\mathcal C$.
\end{lemma}
\begin{proof}
Induct on derivations. Rename local eigenparameters before inserting
proofs or terms. Numerical substitution preserves a formula's root
constructor, so an allowed progress choice remains allowed.
The selected choice does not change under substitution or cut.
For ordinary-assumption cut, replace each use of its constant clause
by the supplied proof. Computation schemas are closed under
capture-free numerical substitution. These facts verify every case.
\end{proof}

\begin{theorem}[Identity-language equivalence]
\label{thm:choice-equivalence}
For every retained clause context $\Gamma$ and native formula $F$,
\[
 \Gamma\vdash_{\mathcal C}F
 \quad\Longleftrightarrow\quad
 \Gamma\vdash_{\mathcal N}F.
\]
The transformations are effective on finite derivations and leave
all numerical and formula syntax unchanged.
\end{theorem}
\begin{proof}
Each new introduction is an instance of the old one after weakening
by the omitted local assumption. For the case choice at an equation,
add the induction instance. For the progress choice, add the zero
instance. Compound case introduction is already the old rule.
Thus every new derivation translates to an old one.

For the reverse direction, transform an old derivation recursively.
Only its atomic DI case needs a construction. Suppose its matrix is
the equation $E(x,y)$ and its transformed premise is
\[
 \pi:\quad\Gamma,E(t,0),E(t,z)\vdash E(t,Sz).
\]
The outer new inference will be case DI with witness $t$. It is
therefore enough to prove $E(t,Sz)$ under
$\Gamma,E(t,0)$.

Rename the eigenvariable in $\pi$ to a fresh $v$. Choose a new
variable $u$ absent from $E(t,y)$ and all retained data. Under
$\Gamma,E(t,0)$, progress DI with witness zero introduces
\[
             \mathsf D_{u,y}E(t,y)
\]
from exactly the renamed premise $\pi[v/z]$.
The bound coordinate $u$ is deliberately unused.

Apply DE to this paired formula. Because $u$ is unused, the opened
clause is exactly
\[
               \langle E(t,0);y.E(t,y)\rangle;
\]
it contains no new existential witness. The retained ordinary
assumption $E(t,0)$ supplies its antecedent. Use at $Sz$ gives
$E(t,Sz)$, and DE closes. Finally the outer case DI introduces
the original $\mathsf D_{x,y}E(x,y)$ with witness $t$.
The original premise proof $\pi$ has been inserted once.

An old compound introduction already is the new case choice.
All other rule cases are unchanged. This completes the induction
on the original derivation.
\end{proof}

The unused inner coordinate is essential. Introducing
$\mathsf D_{x,y}E(x,y)$ internally and then eliminating it would
give an arbitrary fresh witness, not necessarily the original
$t$. The construction instead keeps $t$ as a fixed parameter
of $E(t,y)$ and binds only the unused coordinate $u$.

\begin{proposition}[Linear bound on inference nodes]
\label{prop:choice-node-bound}
Count each computation axiom, Use, DI, DE, and transport as one
inference node, expanding ordinary-assumption use as a computation
of $H$ followed by Use. If an old derivation has $N$ nodes and
$r$ atomic DI nodes, the displayed transformation has exactly
\[
                   N+5r\leq 6N
\]
nodes.
\end{proposition}
\begin{proof}
Each atomic old DI node is replaced by six nodes: outer case DI,
inner progress DI, DE, Use of the opened clause, Use of the
retained base assumption, and the computation of $H$ needed for
that last ordinary-assumption use. Its old premise subtree
occurs once. All other inference nodes are copied once.
\end{proof}

This counts inference-tree nodes. It is not a bound on serialized
JSON bytes, formula sizes, checker runtime, or the number of
consumed packet cells when hypothesis closure is implicit.
In the formula-context Q/D presentation, the corresponding
replacement has outer case DI, inner progress DI, DE, QE,
and two Hyp nodes, giving the same six-node count.

\subsection{Induction and numerical cases separate cleanly}

\begin{lemma}[Contextual equation induction]
\label{lem:choice-ei}
The progress choice, DE, and Use derive
\[
\frac{\Gamma\vdash E(0)\qquad
      \Gamma,E(z)\vdash E(Sz)}
     {\Gamma\vdash E(t)}
\tag{EI}\label{eq:choice-ei}
\]
for every equation $E$, with the usual freshness condition.
\end{lemma}
\begin{proof}
Choose an unused witness coordinate $u$. Progress DI introduces
$\mathsf D_{u,y}E(y)$ from the given step. DE opens the clause
$\langle E(0);y.E(y)\rangle$. Supply its antecedent using the given
base proof and use the clause at $t$. Close DE.
\end{proof}

\begin{lemma}[Cases, equality, and binder congruence]
\label{lem:choice-preinduction}
The case choice derives, for arbitrary native formulas,
\[
\frac{\Gamma\vdash F(0)\qquad\Gamma\vdash F(Sz)}
     {\Gamma\vdash F(t)}
\tag{NC}\label{eq:choice-nc}
\]
with the usual freshness condition. Equality reflexivity,
symmetry, transitivity, numerical congruence, and uniform
matrix interderivability under $\mathsf D$ are derivable
without using the progress choice.
\end{lemma}
\begin{proof}
For cases, use case DI with an unused witness coordinate and
matrix $F(y)$, ignoring its retained base in the successor
premise. DE opens $\langle F(0);y.F(y)\rangle$; the given zero
proof and Use at $t$ finish.

For reflexivity put $r=I(0,t;u.0)$. Use the computation equation
$r=t$ as both transport premises with motive $v=t$ to obtain
$t=t$. The usual explicit equality motives give symmetry,
transitivity, and numerical congruence.

For binder congruence, open the old paired formula at a fresh
witness $a$ and introduce the new one by case DI at $a$.
Convert its retained base to the old base, use the opened
old clause at the successor, and convert back. Ordinary cut
closes the construction, and the reverse conversion gives
the converse. No induction is used.
\end{proof}

The primitive progress choice is still restricted to equations.
Allowing it at arbitrary compound matrices would, by the
argument for \eqref{eq:choice-ei}, supply induction for
arbitrary formulas and change the intended theory.

\subsection{The entire ordinary logical interface}
\label{subsec:choice-logic}

Here is a predecessor-free reconstruction of the ordinary logical
operations. It makes explicit that the case choice supplies the
logical interface independently of progress induction.

Let $C(v,a,b):=I(v,a;u.b)$ with $u$ unused. Computation gives
$C(0,a,b)=a$ and $C(Sv,a,b)=b$. Define the gate on fixed finite
formula syntax by
\[
 G_v(s=t):=(C(v,0,s)=C(v,0,t)),\qquad
 G_v(\mathsf D_{x,y}B):=\mathsf D_{x,y}G_vB,
\]
renaming bound variables away from the control. Equality,
binder congruence, and external structural induction give
\[
 \vdash G_0A,\qquad G_{Sv}A\dashv\vdash A,\qquad
 A\dashv\vdash\mathsf D_{u,y}G_yA
 \quad(u,y\notin\operatorname{FV}(A)).
\tag{Gate and padding}\label{eq:choice-gate}
\]
For padding, case DI proves the forward direction. DE, the
zero-gate theorem, and Use at $1$ prove the reverse direction.

Define a selector between two arbitrary formulas by
\[
 M_v(s=t,u=w):=(C(v,s,u)=C(v,t,w))
\]
and, with aligned fresh paired variables, by
\[
\begin{aligned}
 M_v(\mathsf D_{x,y}A,\mathsf D_{x,y}B)
   &:=\mathsf D_{x,y}M_v(A,B),\\
 M_v(E,\mathsf D_{x,y}B)
   &:=\mathsf D_{x,y}M_v(G_yE,B),\\
 M_v(\mathsf D_{x,y}A,E)
   &:=\mathsf D_{x,y}M_v(A,G_yE),
\end{aligned}
\]
where $E$ is an equation. Every call decreases maximum input
binder depth. Padding and congruence prove
$M_0(A,B)\dashv\vdash A$ and
$M_{Sv}(A,B)\dashv\vdash B$.

Put $\delta(v)=C(v,1,0)$ and $p(z)=I(z,0;u.\delta(u))$.
Only the fixed computations $p(1)=1$ and $p(2)=0$ are needed.
The full logical macro inventory is
\[
\begin{aligned}
 A\Rightarrow_*B&:=\mathsf D_{u,z}M_z(A,B),\\
 A\wedge_*B&:=\mathsf D_{u,z}G_zM_{p(z)}(A,B),\\
 \exists_*x\,A(x)&:=\mathsf D_{x,y}G_yA(x),\\
 A\vee_*B&:=\exists_*x\,M_x(A,B),\\
 \forall_*x\,A(x)&:=
        A(0)\wedge_*\mathsf D_{u,z}G_zA(z),\\
 \bot_*&:=(0=1).
\end{aligned}
\tag{Logical macros}\label{eq:choice-macros}
\]
Newly bound variables are fresh; $u$ is unused.

Implication introduction obtains its assumption from the zero
selector and supplies the consequent at successors, then uses
case DI. Elimination opens its clause and uses it at $1$.
Existential introduction uses its witness and the successor
gate; elimination opens at a fresh witness, supplies the
zero-gate theorem, and uses the clause at $1$.

Disjunction introduction uses witnesses $0$ and $1$.
For elimination, discharge the two branch proofs as
$A\Rightarrow_*F$ and $B\Rightarrow_*F$. NC applied to
$M_z(A,B)\Rightarrow_*F$ gives this implication at every
control: the two endpoint cases follow from the branch
implications. Existential elimination provides $M_a(A,B)$
and the implication gives $F$.

Given $A$ and $B$, selector endpoints and NC give $M_v(A,B)$
at an arbitrary control. Substitute $v=p(Sz)$ and use the
successor gate for conjunction introduction. For projections,
open its zero-gated clause and specialize at $1$ and $2$;
the two fixed computations give $B$ and $A$, respectively.

These rules prove existential numerical coverage:
\[
     (x=0)\vee_*\exists_*v\,(x=Sv).
\tag{Coverage}\label{eq:choice-coverage}
\]
Use NC. Zero takes the left injection of reflexivity, and a
successor takes the right injection with the evident witness.
Coverage gives equality cases even when the retained context
contains $x$; its elimination adds $x=0$ or $x=Sv$ rather
than silently changing that retained context.

Universal introduction follows from a fresh-parameter proof
by substituting zero and successors and forming the two
conjuncts. For elimination, take its $A(0)$ conjunct and
open the paired conjunct. Its base is a zero-gate theorem.
Coverage of the target chooses the zero instance from
$A(0)$ or a successor instance from the opened clause
and gate removal.

Finally, under $0=1$, computation and transport give
$s=C(0,s,t)=C(1,s,t)=t$ for arbitrary numerical terms.
Every equation follows, and external induction on formula
syntax with case DI gives every native formula. This
supplies falsity elimination. No excluded middle for
quantified formulas has entered the construction.

\subsection{Whole-domain realizations and both formula composites}

For a standard first-order reading, let $T_{\rm eq}$ have
intuitionistic equality logic, the same numerical signature
and computation equations, numerical coverage, separation
$(0=1)\to\bot$, and contextual equation induction.
Its interpretation of the native syntax is
\[
 (s=t)^\circ=(s=t),\qquad
 (\mathsf D_{x,y}B)^\circ
   =\exists x\bigl(B^\circ(x,0)\to\forall y B^\circ(x,y)\bigr),
\]
\[
 \langle A;y.B\rangle^\circ
       =A^\circ\to\forall y B^\circ.
\]
The domain is the entire numerical domain and interpreted
equality is actual equality.

Case DI is sound in this realization by numerical coverage.
Progress DI is sound by contextual equation induction under
the conditional base. Use, DE, transport, and computation
have their displayed ordinary meanings. Conversely, the
native macros derive every ordinary logical rule,
\eqref{eq:choice-ei} derives equation induction, and
\eqref{eq:choice-coverage} derives coverage. Separation
translates to the identity implication
$(0=1)\Rightarrow_*(0=1)$.

For the essential paired formula composite,
\[
 \mathsf D_{x,y}B
 \dashv\vdash
 \exists_*x\bigl(B(x,0)\Rightarrow_*\forall_*y\,B(x,y)\bigr).
\tag{Paired composite}\label{eq:choice-paired-composite}
\]
Forwards, DE opens its clause at a fresh $a$; under $B(a,0)$
that clause supplies $B(a,y)$ at a fresh arbitrary $y$.
Introduce the universal, implication, and existential
macros and close DE. Backwards, eliminate the existential
at a fresh $a$. Introduce the paired formula with witness
$a$ by case DI: the retained base, supplied implication,
and universal elimination at $Sz$ give its successor premise.

The ordinary interpretations of the six logical macros
have their ordinary round trips. For conjunction, its true
base and the instances $1,2$ recover both operands;
the converse uses selector cases. The universal has all
successor instances in its paired conjunct and zero in
its first conjunct; coverage gives exactly the ordinary
universal. Implication, disjunction, and existential
quantification use their selector or gate endpoints.
Falsity uses the explicit ordinary separation bridge.
External induction on finite formulas, together with
\eqref{eq:choice-paired-composite}, gives both formula
composites. Clauses correspond in both directions to
$A\Rightarrow_*\forall_*yB$ by the derived logical rules.

The numerical composites are literal identities with
internally provable graph $x=y$. No standard-numeral
restriction, retract, quotient, or equality reflection
assumption is being inserted at this stage. The separately
supplied arithmetic bridge identifies $T_{\rm eq}$ with
full iPRA; the present change preserves that bridge
unchanged because Theorem~\ref{thm:choice-equivalence}
preserves every native sequent.

\subsection{What the theorem does and does not claim}

The inherited Choice-DI interface retains one ordinary assumption
at every paired introduction. It preserves the entire
consequence relation, not just equations or closed
theorems. Consequently the established full-iPRA
interpretations and equiconsistency transfer unchanged.

This result is different from retaining only progress
introduction and strengthening closure elimination.
Here the original elimination rule is retained, and
the case and progress introductions coexist. It is also
different from allowing arbitrary compound progress,
which would supply stronger induction.

A complete-checker comparison must still count the
certificate choice, its equation-root restriction,
substitution, context handling, computation, and the
public input/output interface. The mathematical theorem
alone is not a byte-count or minimality theorem.
The inference-node bound is distinct from source-size
measurements and from serialized certificate size.
The proof is a written universal proof; finite regression
tests do not constitute a Lean/Rocq formalization of it.


\section{Arithmetic strength and both ordinary interpretation composites}
\label{sec:ordinary-bridge}

We spell out the arithmetic bridge to avoid identifying standard numerical
agreement with bi-interpretability. First elaborate every fixed iterator
into an ordinary function symbol. At stage zero the symbols are $0,S$.
Given a fixed earlier-stage term $h(u,\vec p)$, introduce
$i_h(n,a,\vec p)$ with the two equations
\[
i_h(0,a,\vec p)=a,\qquad
i_h(Sn,a,\vec p)=h(i_h(n,a,\vec p),\vec p).
\]
Only finite descriptions occur; the domain contains no code for a
universal evaluator. Formal parameters are distinct when a description
is declared, while actual arguments can coincide.

Let $T_I$ be ordinary intuitionistic first-order logic with actual
equality and genuine falsity, these staged equations, contextual
induction for literal equations, and
\[
\forall x\,(x=0\vee\exists y\,x=Sy),
\qquad (0=S0)\to\bot.
\]
The second sentence is the falsity bridge. It must be stated: positive
numerical equations by themselves admit a singleton structure. Native
explosion from $0=S0$ and a translation of ordinary falsity to that
equation do not authorize silently dropping the bridge from $T_I$.

The native-to-ordinary translation is
\[
(s=t)^\circ=(s=t),\qquad
(\mathsf D_{x,y}B)^\circ=
\exists x\bigl(B^\circ(x,0)\to\forall y B^\circ(x,y)\bigr).
\]
A clause $\langle A;y.B\rangle$ becomes $A^\circ\to\forall yB^\circ$.
For the auxiliary closure used by the compact checkers,
$(\mathsf Q_yB)^\circ=B^\circ(0)\to\forall yB^\circ$.
The interpretation domain is all elements, defined by $x=x$;
equality is actual equality.

Case DI is sound because numerical coverage combines its retained base
with its successor proof. Progress DI is sound by literal-equation
induction under a temporarily assumed base. Use, DE and transport are
ordinary logical inferences under the translations. Conversely the
derived logical operations, equation induction, cases and explosion in
Section~\ref{sec:choice-introduction} recover the rules and axioms of
$T_I$. An auxiliary closure is interderivable with
$\mathsf D_{v,y}B(y)$ for unused $v$: in one direction use case DI with
the closure and its base; in the other eliminate the paired formula
and return its supplied closure. Thus formulas and context clauses have
the stated meanings in both presentations.

\subsection{The arithmetic dependency order}

The scalar update has no recursion-counter argument. We do not conceal
one in an encoding. The inherited internal proofs, reproduced in the
appendices, establish the required arithmetic in this order:
\begin{enumerate}
\item Using equation induction and numerical cases, construct a fixed
scalar predecessor term and prove $P0=0$ and $P(Sn)=n$.
Existential coverage then gives $x=0\vee x=S(Px)$: in the successor
case $x=Sy$, apply $P$ to obtain $Px=y$ and substitute.
\item Construct truncated subtraction and a numerical characteristic
of actual equality. Prove its reflection on the entire original domain.
\item Compile each fixed quantifier-free formula into its characteristic
equation and recover every contextual quantifier-free induction instance.
\item With that induction available, construct pairing and projections
by a bounded square-root search and compile every fixed general
primitive-recursive description into a scalar iterator term.
\end{enumerate}
No pairing or faithful equality detector is assumed to obtain
predecessor. No appeal to standard-model correctness replaces an
internal equality. Every induction fixes its passive parameters;
diagonal substitution occurs only into a previously proved identity.

For a target recursion
\[
f(0,\vec p)=g(\vec p),\qquad
f(Sn,\vec p)=h(n,f(n,\vec p),\vec p),
\]
let $\pi,L,H$ be the proved pairing and projections. The compiler uses
\[
V(n,\vec p)=I\bigl(n,\pi(0,g^*(\vec p));
u.\pi(SL(u),h^*(L(u),H(u),\vec p))\bigr),
\qquad f^*(n,\vec p)=H(V(n,\vec p)).
\]
Equation induction proves $L(V(n,\vec p))=n$; the two projection
identities then prove the target defining equations. The construction
is external recursion on each finite PR description, not an internal
quantification over arbitrary functions.

\begin{theorem}[Whole-domain bi-interpretability and equiconsistency]
The ordinary realization $T_I$ and full first-order iPRA are
bi-interpretable. In both theories $x=y$ defines the isomorphism from
the original domain to its double interpretation. The native Choice-DI
calculus is therefore a faithful proof presentation of a theory
bi-interpretable and equiconsistent with iPRA.
\end{theorem}
\begin{proof}
Interpret an iterator in iPRA by primitive recursion that ignores its
counter argument in the update. iPRA proves its equations, literal
equation induction, coverage and separation. In the reverse direction
use the scalar compiler. The predecessor and equality-reflection proofs
supply the successor axioms and every quantifier-free induction instance.
All ordinary logical rules are preserved.

For two iterators with the same seed and provably equal updates,
equation induction on the equality of their values proves uniqueness at
every counter, with all parameters fixed. The corresponding argument
applies to ordinary primitive recursion. Induct externally on a finite
description: its double translation satisfies its original defining
equations, so internal uniqueness identifies the two functions. This
also proves coherence when substitution changes the staged syntactic
description of an update.

The domain and equality are unchanged; hence the identity graph is
internally total, bijective and preserves every interpreted operation.
The displayed ordinary meanings of the logical macros and the paired
binder give both formula composites by external induction on finite
formula syntax. Clauses are similarly equivalent to their conditional
universal meanings. The explicit falsity bridge preserves contradiction
in both directions. Effective translations of finite proofs give
equiconsistency, without a claim that either theory proves its own
consistency.
\end{proof}


% Adapted with attribution from the pinned unparenthesized one-hole archive.
\section{Internal scalar predecessor bootstrap}
\label{sec:hybrid-bootstrap}

This scalar bootstrap is inherited from the supplied scalar development~\cite{erasure} and is reproduced to keep the scalar checker self-contained. It is not a new contribution of this paper. The proof uses a factor map fixing odd states and a saturated-iteration construction of predecessor.

The available arithmetic terms are $0$, $S$, and
\[
 I(n,a;u.h),\qquad
 I(0,a;u.h)=a,\qquad
 I(Sn,a;u.h)=h[I(n,a;u.h)/u].
\]
The update binds only $u$.  All other variables in it are fixed numerical
parameters.  We use equality reasoning, the already reconstructed
zero/successor cases, and induction on a single literal equation.
Every induction below keeps its other parameters fixed.  Predecessor,
successor injectivity, pairing, and induction on arbitrary
quantifier-free formulas are not assumptions of the argument.

\subsection{Elementary terms and two finite cycles}

Write $1=S0$ and use larger numerals as successor abbreviations.  Define
\[
 C(n,a,b)=I(n,a;u.b),\qquad
 Z(n)=C(n,1,0),\qquad
 a+b=I(b,a;u.Su).
\]
The update variable of $C$ is absent from its body.  Consequently
\[
 C(0,a,b)=a,\qquad C(Sn,a,b)=b.
\]
Numerical cases show that $Z(n)$ is $0$ or $1$.
The usual commutative-monoid equations for addition follow by literal
equation induction.  For example, $0+b=b$ and $Sa+b=S(a+b)$ are
proved by induction on $b$; these yield commutativity by a further
equation induction, and associativity follows by induction on the last
summand.  We use these equations to reorder finite sums.  The notation
$2a$ means $a+a$ and $3a$ means $a+a+a$; no multiplication operation
is needed for this bootstrap.

Put
\[
 N(n)=I(n,0;u.Z(u)),\qquad
 j(u)=Z(u)+2N(u),\qquad
 M(n)=I(n,0;u.j(u)).
\]
The range of $N$ is contained in $\{0,1\}$: its zero value is $0$,
and its successor value is a value of $Z$.  This uses numerical cases
on its counter, not induction on a disjunction.  Hence
\begin{equation}
 N(n)+Z(N(n))=1,\qquad
 Z(Z(N(n)))=N(n),\qquad N(n+2)=N(n).
 \label{eq:hybrid-parity}
\end{equation}
Here the last equation follows from two successor computations and
the preceding equation.  Literal-equation inductions give
\begin{equation}
 N(2n)=0,\qquad N(2n+1)=1,\qquad N(3n)=N(n).
 \label{eq:hybrid-parity-scaling}
\end{equation}
For the first equation the successor step adds two to the argument of
$N$; the second follows by one computation.  In the third, increasing
$n$ adds three to the left argument and therefore applies $Z$ once to
both sides of the induction equation.

Finite computations now give
\[
 j(0)=1,\qquad j(1)=2,\qquad j(2)=0.
\]
The proof below needs just these three values and
$M(Sn)=j(M(n))$; it does not assume a division algorithm or a theorem
characterizing $M$ as remainder modulo three.

\subsection{A ramp with a global scaling identity}

Define
\[
 G(n)=I(n,0;u.u+M(u)+1).
\]
The central invariant is
\begin{equation}
 M(G(n))=N(n).
 \label{eq:hybrid-ramp-invariant}
\end{equation}
Its zero case is immediate.  For its successor step, retain the
single induction equation and split the Boolean value $N(n)$.
If it is $0$, the next $G$ state is $G(n)+1$; its $M$ value is
$j(0)=1$, equal to $N(Sn)$.  If it is $1$, the next $G$ state
is $G(n)+2$; two computations of $M$ change its value from $1$
to $2$ and then to $0$, again equal to $N(Sn)$.  This proves
the invariant by literal-equation induction.

It follows that
\[
 G(Sn)=G(n)+N(n)+1,\qquad G(n+2)=G(n)+3.
\]
Equation induction therefore yields
\begin{equation}
 G(2n)=3n.
 \label{eq:hybrid-ramp-double}
\end{equation}
A separate literal-equation induction proves
\begin{equation}
 G(3n)=3G(n)+N(n).
 \label{eq:hybrid-ramp-triple}
\end{equation}
For completeness, its successor calculation is
\begin{align*}
 G(3Sn)
 &=G(3n)+N(3n)+4
   =3G(n)+2N(n)+4,\\
 3G(Sn)+N(Sn)
 &=3G(n)+3N(n)+3+Z(N(n)).
\end{align*}
The expressions agree in both Boolean cases for $N(n)$.
The base is $0=0$.  In particular,
\eqref{eq:hybrid-ramp-triple} was proved internally; it was not
inferred from an intended floor-function interpretation of $G$.

Set
\[
 H(u)=C(N(u),G(u),u).
\]
The crucial choice is to retain $u$ when its parity is $1$.
The preceding equations give
\begin{equation}
 H(2u)=3u,\qquad H(1)=1,\qquad H(3u)=3H(u).
 \label{eq:hybrid-factor-map}
\end{equation}
To check the last identity, use $N(3u)=N(u)$ and split the Boolean
$N(u)$.  At $0$, both sides are $3G(u)$ by
\eqref{eq:hybrid-ramp-triple}; at $1$, both are $3u$.
This argument does not require that an arbitrary number has already
been represented as $2v$ or $2v+1$.

\subsection{A comparator from scalar iteration}

For every fixed update $h$, equation induction on $n$, keeping $m,a$
and all parameters fixed, proves the semigroup identity
\[
 I(m+n,a;h)=I(n,I(m,a;h);h).
\]
The base is the zero computation and the step applies $h$ to the
induction equation. This identity will also be used in the later
pairing construction.


For a fixed update $h$, equation induction on $y$, with its initializer
$a$ and all parameters of $h$ fixed, proves the rotation identity
\begin{equation}
 I(Sy,a;h)=I(y,h(a);h).
 \label{eq:hybrid-rotation}
\end{equation}
Both sides have base $h(a)$, and their successor equations apply the
same $h$ to the previous equality.  Here and below $I(y,a;h)$ is
an abbreviation for one fixed finite scalar term, not quantification
over functions.

By \eqref{eq:hybrid-factor-map}, another equation induction gives
\begin{equation}
 I(y,3a;H)=3I(y,a;H).
 \label{eq:hybrid-iterator-scaling}
\end{equation}
Its base is $3a=3a$.  In the step the fixed-parameter induction
equation rewrites the new left side as
\[
 H(3I(y,a;H))=3H(I(y,a;H)).
\]
The initializer parameter stays $a$ throughout this induction.
Equation induction also gives $I(y,1;H)=1$, since $H(1)=1$.

Define
\[
 E(x)=I(x,1;u.2u),\qquad
 L(x,y)=N\bigl(I(y,E(x);H)\bigr).
\]
Then the following entire recursion specification is provable:
\begin{equation}
 L(0,y)=1,\qquad L(Sx,0)=0,\qquad L(Sx,Sy)=L(x,y).
 \label{eq:hybrid-comparator}
\end{equation}
The first equation uses $E(0)=1$ and the fixed-point identity for
$H$.  The second uses $E(Sx)=2E(x)$ and $N(2a)=0$.
For the third,
\begin{align*}
 L(Sx,Sy)
 &=N\bigl(I(Sy,2E(x);H)\bigr)\\
 &=N\bigl(I(y,H(2E(x));H)\bigr)\\
 &=N\bigl(I(y,3E(x);H)\bigr)\\
 &=N\bigl(3I(y,E(x);H)\bigr)\\
 &=N\bigl(I(y,E(x);H)\bigr)=L(x,y).
\end{align*}
The successive steps use rotation, $H(2a)=3a$, iterator scaling,
and $N(3a)=N(a)$.  Thus $L$ has the comparison equations at all
arguments internally.  Its standard-model meaning as the indicator
of $x\le y$ is motivation, not a premise of the following argument.

\subsection{Saturated iteration gives predecessor}

We now need only the three equations
\eqref{eq:hybrid-comparator}.  Define
\[
 B_m(u)=C(L(Su,m),u,Su),\qquad
 W(t,m)=I(t,0;u.B_m(u)).
\]
The subscript in $B_m$ records a fixed parameter; $B_m$ is an open
term abbreviation, not a function variable.

Numerical cases on $e$ prove the selector equation
\[
 C(e,Sa,Sb)=S C(e,a,b).
\]
Equivalently one may use equation induction on $e$, whose induction
hypothesis is unused.  Therefore
\begin{align}
 B_{Sm}(Su)
 &=C(L(SSu,Sm),Su,SSu)\notag\\
 &=C(L(Su,m),Su,SSu)=S B_m(u),
 \label{eq:hybrid-conjugacy}\\
 B_{Sm}(0)&=C(L(S0,Sm),0,S0)=C(L(0,m),0,S0)=S0.
 \label{eq:hybrid-conjugacy-base}
\end{align}

Literal-equation induction on $t$, with $m$ fixed, proves
\begin{equation}
 W(St,Sm)=S W(t,m).
 \label{eq:hybrid-saturated-shift}
\end{equation}
The base is \eqref{eq:hybrid-conjugacy-base}.  The successor step is
\begin{align*}
 W(SSt,Sm)
 &=B_{Sm}(W(St,Sm))\\
 &=B_{Sm}(S W(t,m))\\
 &=S B_m(W(t,m))=S W(St,m).
\end{align*}
Only the second line uses the current induction hypothesis, at the
same fixed $m$.

Two subsequent equation inductions give
\begin{equation}
 L(Sn,n)=0,\qquad W(n,n)=n.
 \label{eq:hybrid-diagonals}
\end{equation}
The first uses the comparator recurrence.  The second uses
\eqref{eq:hybrid-saturated-shift} in its successor step.  The
substitution $t=m=n$ is made into that already proved universally
parameterized equation; no parameter of an active induction hypothesis
is changed.  Consequently
\begin{equation}
 W(Sn,n)=B_n(n)=C(L(Sn,n),n,Sn)=n.
 \label{eq:hybrid-extra-step}
\end{equation}

Finally define the finite scalar term
\begin{equation}
 P(m)=I\bigl(m,0;u.C(L(SSu,m),u,Su)\bigr).
 \label{eq:hybrid-predecessor-term}
\end{equation}
Its zero equation is immediate.  To justify its successor equation
without assuming a separate extensionality rule for iterator bodies,
put
\[
 V(t,m)=I\bigl(t,0;u.C(L(SSu,Sm),u,Su)\bigr).
\]
The comparator recurrence identifies this update pointwise with
$B_m(u)$.  Equation induction on $t$, with $m$ fixed, gives
$V(t,m)=W(t,m)$: their base values are $0$, and the successor step
uses the pointwise update equation and then the induction equation.
Now \eqref{eq:hybrid-extra-step} yields
\[
 \boxed{\qquad P(0)=0,\qquad P(Sm)=V(Sm,m)=W(Sm,m)=m.\qquad}
\]

\subsection{Dependency boundary}

All terms in this section are built from $0,S,I$ by finite substitution.
All inductions have a single equation as their invariant.  Numerical
cases and equality transport are the already reconstructed logical
resources.  Successor injectivity may now be obtained by applying $P$
to an assumed equation $Sx=Sy$.  Equality tests, reflection of
quantifier-free formulas, pairing, and the recovery of general
primitive recursion can then be developed in the established order.
They were not used in order to manufacture predecessor.

This argument is a written internal derivation in the stated theory.
The finite implementation checks do not constitute a proof-assistant
formalization of this entire metatheorem.


% Standalone section.  Requires only amsmath and amssymb in the parent document.
% No theorem environments, macros, or bibliography keys are assumed.
% All arithmetic claims in this section are proved here.

\section{From literal-equation induction to full intuitionistic PRA}
\label{sec:eic-completion}\label{app:arithmetic}

This section supplies an unconditional completion theorem. It is adapted
with attribution from the current forward-left draft~\cite{forward}. In the
scalar application used here, every function occurring in the proof is
already a scalar term: predecessor is supplied by
Appendix~\ref{sec:hybrid-bootstrap}, the selector is
$C(n,a,b)=I(n,a;u.b)$, addition is $x+y=I(y,x;u.Su)$, and
truncated subtraction is $d(x,y)=I(y,x;u.Pu)$. All other functions below
are finite compositions of these. Thus this proof does not presuppose
that arbitrary primitive recursion has already been compiled. The QFI
conclusion first applies to the scalar language; the general PR compiler
comes afterward.  Literal-equation
induction, together with the primitive-recursive defining equations,
predecessor coverage, and separation of $0$ and $1$, proves induction for every
quantifier-free formula.  In particular, the proof does not assume a faithful
numerical equality test, stability or decidability of equality, successor
uniqueness, or any commutation law for an iterator.  We construct the equality
test and prove its reflection property for arbitrary elements of the
numerical domain.

\subsection{The exact hypotheses}

Work in ordinary intuitionistic first-order logic with equality and genuine
falsity $\bot$.  The language has one numerical sort and a symbol for each
fixed, finite primitive-recursive description.  Distinct descriptions are
not identified merely because they compute the same standard function.
Let $\mathsf{PRE}$ contain the defining equations for zero, successor,
projections, composition, and primitive recursion.  In particular, choose
descriptions $P$ and $C$ with equations
\begin{align}
 P(0)&=0, & P(Sx)&=x, \label{eq:eic-predecessor}\\
 C(0,a,b)&=a, & C(Sx,a,b)&=b. \label{eq:eic-selector}
\end{align}
Write $1=S0$.  In addition to $\mathsf{PRE}$, assume
\begin{align}
 \mathsf{Cov}:&\qquad
       \forall x\,\bigl(x=0\lor x=S(Px)\bigr),
       \label{eq:eic-coverage}\\
 \mathsf{Sep}:&\qquad 0=1\longrightarrow\bot,
       \label{eq:eic-separation}
\end{align}
and induction for each \emph{literal equation} $E(n,\vec p)$:
\begin{equation}
 \forall\vec p\,
 \left[
   \left(E(0,\vec p)\land
       \forall n\bigl(E(n,\vec p)\longrightarrow E(Sn,\vec p)\bigr)
   \right)
   \longrightarrow\forall n\,E(n,\vec p)
 \right].
 \label{eq:eic-induction}
\end{equation}
Here the schematic expression $E$ is a single equality of numerical
terms; the displayed universal quantifiers range only over numbers.
Every description and every instance of the schema is finite external
syntax.

The axiom formulation \eqref{eq:eic-induction} and the usual contextual
induction rule are interchangeable.  For example, if
\[
 \Gamma\vdash E(0,\vec p),
 \qquad
 \Gamma,E(z,\vec p)\vdash E(Sz,\vec p),
\]
where $z$ is absent from $\Gamma$ and the passive parameters, discharge the
second premise's equation and universally introduce $z$.  The axiom then
gives $\Gamma\vdash\forall n E(n,\vec p)$, which can be instantiated at any
target term.  Conversely, assume the antecedent of the axiom and apply the
contextual rule with fresh variables.  This uses the ordinary intuitionistic
logical rules and permits arbitrary, including quantified, formulas in
$\Gamma$.

\subsection{Constructor separation and two elementary tests}

Define the sign description and the countdown description by
\begin{equation}
 \operatorname{sg}(u)=C(u,0,1),
 \qquad
 K(u)=C(u,1,Pu).
 \label{eq:eic-sign-countdown}
\end{equation}
Their defining equations imply
\begin{align}
 \operatorname{sg}(0)&=0,
 &\operatorname{sg}(Sx)&=1,\label{eq:eic-sign-computation}\\
 K(0)&=1,
 &K(Sx)&=x.\label{eq:eic-countdown-computation}
\end{align}

\paragraph{Constructor laws.}
Successor is injective: applying $P$ to $Sx=Sy$ gives $x=y$ by
\eqref{eq:eic-predecessor}.  Successor never equals zero.  Indeed,
$Sx=0$ gives $1=0$ by applying $\operatorname{sg}$ and using
\eqref{eq:eic-sign-computation}; symmetry and
\eqref{eq:eic-separation} give $\bot$.  Consequently $1\ne0$, where
$A\ne B$ abbreviates $A=B\to\bot$.  None of these arguments uses
induction.

\paragraph{The sign test.}
For every $u$,
\begin{align}
 &\operatorname{sg}(u)=0\lor\operatorname{sg}(u)=1,
       \label{eq:eic-sign-range}\\
 &\operatorname{sg}(u)=0\quad\longleftrightarrow\quad u=0,
       \label{eq:eic-sign-zero}\\
 &\operatorname{sg}(u)=1\quad\longleftrightarrow\quad u\ne0.
       \label{eq:eic-sign-one}
\end{align}
To prove the range assertion, split \eqref{eq:eic-coverage}.  The branch
$u=0$ gives sign $0$, and the branch $u=S(Pu)$ gives sign $1$.
The implication $u=0\to\operatorname{sg}(u)=0$ is computation.
For its converse, the successor branch of coverage would make
$\operatorname{sg}(u)=0$ imply $1=0$, a contradiction; in that branch
explosion gives $u=0$ as well.  Similarly, if
$\operatorname{sg}(u)=1$ and $u=0$, computation gives $0=1$.
Conversely, $u\ne0$ excludes the zero branch of coverage, and the
successor branch computes the sign as $1$.

Coverage and \eqref{eq:eic-selector} also give, for all $u,v$,
\begin{equation}
 C(u,v,v)=v.
 \label{eq:eic-selector-idempotence}
\end{equation}
Both coverage branches compute to $v$, so this identity again needs no
induction.

\paragraph{The countdown test and the notation $W$.}
For every $u$,
\begin{equation}
 K(u)=0\quad\longleftrightarrow\quad u=1.
 \label{eq:eic-countdown-zero}
\end{equation}
If $u=1=S0$, \eqref{eq:eic-countdown-computation} gives $K(u)=0$.
Conversely, assume $K(u)=0$ and split coverage.  If $u=0$, then
$K(u)=1$, contradicting separation, and explosion gives $u=1$.
If $u=S(Pu)$, computation gives $K(u)=Pu$; hence $Pu=0$ and
$u=S0=1$.

If a presentation uses a description named $W$ with defining equations
\[
 W(0)=1,\qquad W(Sx)=x,
\]
then it proves $K(u)=W(u)$ for every $u$ already from coverage.
In the zero branch both sides are $1$; in the branch $u=S(Pu)$ both
sides are $Pu$.  Thus $K$ and $W$ are provably equal on the whole domain,
even if they are distinct descriptions.  This equality does not appeal to
standard-model agreement or to uniqueness of an arbitrary recurrence.

\subsection{Four elementary uses of equation induction}

Choose descriptions for addition and truncated subtraction, with recursion
on the second argument:
\begin{align}
 x+0&=x, &x+Sy&=S(x+y),\label{eq:eic-addition}\\
 d(x,0)&=x, &d(x,Sy)&=P(d(x,y)).\label{eq:eic-subtraction}
\end{align}
Using these equations and \eqref{eq:eic-induction}, we now derive
\begin{align}
 0+y&=y,\label{eq:eic-left-zero}\\
 (Sx)+y&=S(x+y),\label{eq:eic-left-successor}\\
 d(Sx,Sy)&=d(x,y),\label{eq:eic-cancel-successors}\\
 d(Sa,a)&=1.\label{eq:eic-diagonal}
\end{align}
Each induction below has a single equation as its induction formula.

For \eqref{eq:eic-left-zero}, induct on $y$.  The base is $0+0=0$.
The step is $0+Sy=S(0+y)=Sy$.

For \eqref{eq:eic-left-successor}, hold $x$ fixed and induct on $y$.
At zero both sides are $Sx$.  At a successor the induction hypothesis
gives
\[
 (Sx)+Sy=S\bigl((Sx)+y\bigr)
       =S\bigl(S(x+y)\bigr)=S(x+Sy).
\]
Combining \eqref{eq:eic-left-successor} with the defining equation for
addition gives the useful rearrangement
\begin{equation}
 (Sy)+r=y+Sr.
 \label{eq:eic-shift-successor}
\end{equation}
Indeed, both sides are $S(y+r)$.

For \eqref{eq:eic-cancel-successors}, hold $x$ fixed and induct on $y$.
The base computes as
\[
 d(Sx,S0)=P(d(Sx,0))=P(Sx)=x=d(x,0).
\]
Applying $P$ to the induction equation gives the step:
\[
 d(Sx,S(Sy))=P(d(Sx,Sy))=P(d(x,y))=d(x,Sy).
\]

Finally, induct on $a$ for \eqref{eq:eic-diagonal}.  The base is
$d(S0,0)=S0=1$.  At a successor,
\[
 d(S(Sa),Sa)=d(Sa,a)=1,
\]
using \eqref{eq:eic-cancel-successors} and the induction hypothesis.

\subsection{An equation invariant that reflects actual equality}

The remaining arithmetic ingredient is the following equation, valid for
all $x,y$:
\begin{equation}
 C\bigl(\operatorname{sg}(d(x,y)),\ x,\ y+d(x,y)\bigr)=x.
 \label{eq:eic-residual-invariant}
\end{equation}
Its meaning is that a nonzero subtraction residual reconstructs the
original argument by addition.  The selector expresses this assertion as
one literal equation, so equation induction applies directly.

Hold $x$ fixed and induct on $y$.  At $y=0$, subtraction gives residual
$x$, and \eqref{eq:eic-left-zero} gives $0+x=x$.  The desired equation
is consequently $C(\operatorname{sg}(x),x,x)=x$, which is
\eqref{eq:eic-selector-idempotence}.

For the step, assume \eqref{eq:eic-residual-invariant} at $y$ and
abbreviate
\[
 r=d(x,y),\qquad r'=Pr=d(x,Sy).
\]
These are local term abbreviations, not additional function symbols.
Split coverage for $r'$.

If $r'=0$, then $\operatorname{sg}(r')=0$, and the successor instance
of \eqref{eq:eic-residual-invariant} reduces immediately to $x=x$.

In the other branch, $r'=S(Pr')$, so constructor separation gives
$r'\ne0$.  Split coverage for $r$.  Its zero branch would imply
$r'=P0=0$, a contradiction, from which the desired step equation follows.
In its successor branch,
\[
 r=S(Pr)=Sr',\qquad \operatorname{sg}(r)=1.
\]
The induction equation therefore computes to $y+r=x$.  Meanwhile,
\eqref{eq:eic-shift-successor} gives
\[
 (Sy)+r'=y+Sr'=y+r=x.
\]
Since $r'$ is a successor, $\operatorname{sg}(r')=1$, so this is
exactly the branch selected in the successor instance of
\eqref{eq:eic-residual-invariant}.  This completes the equation induction.
All case distinctions came from coverage; no induction on an implication
or a disjunction was used.

\paragraph{Reflection at residual $1$.}
We have, for arbitrary $a,b$,
\begin{equation}
 d(Sa,b)=1\quad\longleftrightarrow\quad a=b.
 \label{eq:eic-residual-reflection}
\end{equation}
For the forward direction, substitute $x=Sa$, $y=b$ into
\eqref{eq:eic-residual-invariant}.  If $d(Sa,b)=1$, the sign is $1$
and the selector returns its second branch, giving
\[
 b+1=Sa.
\]
The defining equation for addition gives $b+1=S(b+0)=Sb$.  Thus
$Sb=Sa$, and applying $P$ gives $b=a$, hence $a=b$.
Conversely, $a=b$ and \eqref{eq:eic-diagonal} give
$d(Sa,b)=d(Sa,a)=1$.

Define the equality characteristic by the fixed description
\begin{equation}
 e(a,b)=\operatorname{sg}\bigl(K(d(Sa,b))\bigr).
 \label{eq:eic-equality-characteristic}
\end{equation}
The sign, countdown, and residual reflection lemmas give
\begin{align}
 e(a,b)=0&\quad\longleftrightarrow\quad a=b,
       \label{eq:eic-equality-reflection}\\
 e(a,b)=0&\ \lor\ e(a,b)=1.
       \label{eq:eic-equality-range}
\end{align}
In detail, $e(a,b)=0$ is equivalent to $K(d(Sa,b))=0$ by
\eqref{eq:eic-sign-zero}, which is equivalent to $d(Sa,b)=1$ by
\eqref{eq:eic-countdown-zero}, and then to $a=b$ by
\eqref{eq:eic-residual-reflection}.  The range assertion is
\eqref{eq:eic-sign-range} applied to $K(d(Sa,b))$.

The reflection property concerns the logic's actual equality throughout
the domain.  It was not assumed as part of a numerical comparison symbol,
and its proof did not restrict $a$ or $b$ to standard numerals.
It also yields decidable equality.  In the $e(a,b)=0$ branch,
\eqref{eq:eic-equality-reflection} gives $a=b$.  In the $e(a,b)=1$
branch, an assumption $a=b$ would give $e(a,b)=0$ and therefore $0=1$,
contradicting separation.  Hence
\begin{equation}
 \forall a\,\forall b\,\bigl(a=b\lor(a=b\to\bot)\bigr)
 \label{eq:eic-equality-decidable}
\end{equation}
is a theorem, not an extra hypothesis.

\subsection{Compilation of all quantifier-free formulas}

For each fixed quantifier-free formula $A(\vec p)$, construct a term
$b_A(\vec p)$ by external recursion on its finite syntax.  Take equality,
$\bot$, $\land$, $\lor$, and $\to$ as the primitive formula constructors;
negation is implication into $\bot$.  Set
\begin{align}
 b_{s=t}&=e(s,t), & b_{\bot}&=1,\label{eq:eic-boolean-atoms}\\
 b_{A\land D}&=C(b_A,b_D,1),\label{eq:eic-boolean-and}\\
 b_{A\lor D}&=C(b_A,0,b_D),\label{eq:eic-boolean-or}\\
 b_{A\to D}&=C(b_A,b_D,0).\label{eq:eic-boolean-implies}
\end{align}
Every term here is a finite composition of fixed primitive-recursive
descriptions; if desired, projections and composition give it a single
description symbol.  No internal quantification over formulas or programs
is involved.

An external induction on the syntax of $A$ proves, within the arithmetic
theory,
\begin{equation}
 \bigl(b_A=0\lor b_A=1\bigr)
 \quad\text{and}\quad
 \bigl(b_A=0\longleftrightarrow A\bigr).
 \label{eq:eic-boolean-correctness}
\end{equation}
Here are the full connective checks.  In every case $b_A=1$ implies
$\neg A$: if $A$ held, the already established implication
$A\to b_A=0$ would contradict separation.

\paragraph{Atomic formulas and falsity.}
For an equation, both claims are
\eqref{eq:eic-equality-reflection}--\eqref{eq:eic-equality-range}.
For falsity the range assertion is $1=0\lor1=1$, with its right disjunct
proved by reflexivity.  Also $1=0\to\bot$ follows from separation and
symmetry, while $\bot\to1=0$ follows by explosion.  Thus $1=0$ is
equivalent to $\bot$, as required.

\paragraph{Conjunction.}
Split $b_A=0\lor b_A=1$.  In the zero branch, $A$ holds and
$b_{A\land D}=b_D$.  Its range follows from the range of $b_D$,
and its being zero is equivalent to $D$, hence to $A\land D$ under
the known $A$.  In the one branch, $\neg A$ holds and
$b_{A\land D}=1$.  Its range is immediate.  If it were zero,
separation would give $\bot$, hence $A\land D$; conversely an
assumption $A\land D$ contradicts $\neg A$, and explosion gives
$b_{A\land D}=0$.

\paragraph{Disjunction.}
In the branch $b_A=0$, the formula $A$ holds and
$b_{A\lor D}=0$.  The range is immediate, and both $b_{A\lor D}=0$
and $A\lor D$ hold, proving their equivalence.  In the branch
$b_A=1$, we have $\neg A$ and $b_{A\lor D}=b_D$.  Its range follows
from the range of $b_D$.  If its value is zero, then $D$ and hence
$A\lor D$.  Conversely, from $A\lor D$, the $A$ alternative
contradicts $\neg A$ and yields $b_D=0$ by explosion, while the
$D$ alternative yields $b_D=0$ directly.  This proves the required
equivalence in both alternatives.

\paragraph{Implication.}
In the branch $b_A=0$, the formula $A$ holds and
$b_{A\to D}=b_D$.  Its range follows from the range of $b_D$.
If its value is zero, then $D$, which gives $A\to D$.  Conversely,
$A\to D$ and the known $A$ give $D$, hence $b_D=0$.
In the branch $b_A=1$, we have $\neg A$ and
$b_{A\to D}=0$.  Its range is immediate.  To prove $A\to D$,
assume $A$; this contradicts $\neg A$, and explosion gives $D$.
Thus both $b_{A\to D}=0$ and $A\to D$ hold in this branch, and
their equivalence follows.  This argument uses the proved two-value
range of the finite characteristic.  It does not use excluded middle
for arbitrary formulas or change the surrounding intuitionistic logic.

This completes the external syntactic induction proving
\eqref{eq:eic-boolean-correctness} for every quantifier-free formula.

\subsection{Recovery of induction and the full theory}

Let $A(n,\vec p)$ be any fixed quantifier-free formula.  Suppose an
arbitrary retained context $\Gamma$ proves its base and its induction
step:
\[
 \Gamma\vdash A(0,\vec p),
 \qquad
 \Gamma,A(z,\vec p)\vdash A(Sz,\vec p),
\]
with the usual freshness condition on $z$.  By
\eqref{eq:eic-boolean-correctness}, the base gives
$\Gamma\vdash b_A(0,\vec p)=0$.  Under the additional assumption
$b_A(z,\vec p)=0$, the same equivalence gives $A(z,\vec p)$,
the given step gives $A(Sz,\vec p)$, and the equivalence gives
$b_A(Sz,\vec p)=0$.  Apply literal-equation induction to
$b_A(n,\vec p)=0$ and then apply its reflection implication.
The result is $\Gamma\vdash A(t,\vec p)$ for every target term $t$.
Thus the theory proves the full contextual quantifier-free induction
rule, including in arbitrary quantified contexts.

\paragraph{Completion theorem.}
In the common language of fixed primitive-recursive descriptions, the
theory
\[
 \mathsf{PRE}+\mathsf{Cov}+\mathsf{Sep}+\mathsf{EI}
\]
has exactly the same sequents as ordinary first-order intuitionistic
primitive recursive arithmetic with all quantifier-free induction
instances.  The forward simulation follows from the recovered
quantifier-free induction, successor injectivity, and successor separation;
all primitive-recursive defining equations and all logical rules are
shared.  For the reverse simulation, literal-equation induction is an
instance of quantifier-free induction, and successor separation gives
$\mathsf{Sep}$.  Quantifier-free induction proves $\mathsf{Cov}$:
at zero choose its left disjunct, and at a successor use
$Sx=S(P(Sx))$, which follows from $P(Sx)=x$, for its right disjunct.

Both interpretations therefore use the full domain, defined by $x=x$,
and preserve every function symbol, equality, logical connective, and
quantifier.  Both composites have the identity isomorphism, with graph
$x=y$.  The theories are consequently identity-domain bi-interpretable
and equiconsistent.  This completion theorem is unconditional once the
literal-equation induction rule has been obtained; it has no dependency
on a classical successor-uniqueness reduction.


\section{Pairing, primitive recursion, and internal numerical composites}
\label{sec:iterator-coding}
This appendix is adapted with attribution from the progress-closure
arithmetic development~\cite{progress}. Its prerequisites have now been
proved: scalar predecessor, arithmetic reflection, and full quantifier-free
induction. Pairing is a derived proof tool; it is not part of any checker.

\subsection{Pairing from a bounded square-root search}

Once QFI is available, a familiar square-based injection gives a
particularly direct compiler. Write $a\le b$ for $d(a,b)=0$.
The usual linear-order and semiring identities follow by fixed
quantifier-free inductions from the successor, addition, and monus
equations. Multiplication itself is the scalar term
\[
                        a\cdot b=I(b,0;u.u+a).
\]
Only primitive-recursive arithmetic is involved in the following
bounded root search; there is no minimization over an unbounded domain.

For a fixed parameter $n$, put
\[
 U_n(v)=C(d((Sv)^2,n),Sv,v),\qquad
 R_n(i)=I(i,0;v.U_n(v)),\qquad q(n)=R_n(n).
\]
The update increments $v$ exactly when $(v+1)^2\le n$, and
otherwise retains it. QFI in $i$, with $n$ fixed, proves
\[
 R_n(i)\le i,\qquad R_n(i)^2\le n,\qquad
 R_n(i)=i\ \vee\ n<(S(R_n(i)))^2.              \tag{Root invariant}
\]
The base has $R_n(0)=0$. If the update increments, its test
contradicts the strict-upper-bound alternative in the third invariant.
Thus the old state equals $i$, and the new state equals $Si$;
its square bound is the increment test. If the update does not increment, the
strict upper square bound holds and the state remains fixed.
An already established strict upper square bound forces every
subsequent update to remain fixed. These are exactly the cases
needed for the three displayed invariants.

At $i=n$, either the strict upper square bound already holds or
$R_n(n)=n$, in which case $n<(n+1)^2$ supplies it. Thus
\[
                        q(n)^2\le n<(q(n)+1)^2.           \tag{Root}
\]
Monotonicity of squares shows uniqueness of a number satisfying
Root. All formulas just used are quantifier free in the scalar
language, with fixed $n$ during the induction.

Now define
\[
 \pi(a,b)=(a+b)^2+a,\qquad
 L(n)=d(n,q(n)^2),\qquad H(n)=d(q(n),L(n)).
\]
For $k=a+b$, one has $k^2\le k^2+a<(k+1)^2$ since $a\le k$.
Root uniqueness gives $q(\pi(a,b))=a+b$, and therefore
\[
                    L(\pi(a,b))=a,\qquad H(\pi(a,b))=b.
\]
Surjectivity of $\pi$ is unnecessary. The original numerical domain
is never replaced by the range of this coding; pairing is used only
to store the finite state of a fixed primitive recursion. This
standard square pairing is a convenient reconstruction here, not
a claimed new arithmetic basis.

\subsection{Compiling each primitive recursion}

Suppose the fixed target description is
\[
 f(0,\vec p)=g(\vec p),\qquad
 f(Sn,\vec p)=h(n,f(n,\vec p),\vec p).
\]
Assume by external structural induction that the earlier $g,h$
already have scalar translations $g^*,h^*$. Define
\[
 W(u,\vec p)=\pi(SL(u),h^*(L(u),H(u),\vec p)),
\]
\[
 V(n,\vec p)=I(n,\pi(0,g^*(\vec p));u.W(u,\vec p)),
 \qquad f^*(n,\vec p)=H(V(n,\vec p)).
\]
The update has only its state variable bound, with every $p_i$
fixed. Equation induction proves $L(V(n,\vec p))=n$.
The two projection equations and scalar computation then prove
both defining equations for $f^*$ internally. Initial functions
and composition translate directly. This completes the translation
of every fixed finite PR description without adding a universal
evaluator.

\subsection{Why both numerical round trips are internal}

Two fixed scalar iterators with the same seed and provably equal
updates are internally equal at all counters: use literal-equation
induction comparing their values. Their base agrees, and their
successor equations plus congruence prove the step. The same
uniqueness argument works for ordinary primitive recursion.

For an ordinary PR description, induct externally on its finite
construction. The double translation satisfies its original
defining equations, so internal equation induction identifies it
with the original operation. For a native scalar description,
earlier translated updates agree by the external induction
hypothesis; scalar uniqueness identifies the double-translated
iterator. These proofs hold at arbitrary internal counters.

Thus the domain relation $x=y$ is provably a total bijection
preserving every fixed numerical operation in both composites.
The formula constructions and round trips in
Section~\ref{sec:choice-introduction} give the corresponding formula
composites. Both interpretations preserve falsity, taking native
$0=1$ to an ordinary contradiction and ordinary $\bot$ to native
$0=1$. Finite proof translations therefore establish equiconsistency.


\section{The complete new checker}
\label{sec:listing}

This is the full standalone program. Its source-size coordinates include
all the code below. The shorter normalized rendering differs only in
formatting.
\begin{lstlisting}[language=Python]
import json
import sys


def require(condition):
    if not condition:
        raise ValueError('Invalid certificate')


def term(a, n):
    match a:
        case ['v', i]:
            require(type(i) is int and 0 <= i < n)
        case ['z']:
            pass
        case ['s', b]:
            term(b, n)
        case ['i', c, b, h]:
            term(c, n)
            term(b, n)
            term(h, n + 1)
        case _:
            raise ValueError('Invalid term')


def formula(a, n):
    match a:
        case ['=', s, t]:
            term(s, n)
            term(t, n)
        case ['q', b]:
            formula(b, n + 1)
        case ['d', b]:
            formula(b, n + 2)
        case _:
            raise ValueError('Invalid formula')


def walk(a, f, d=0):
    if a[0] == 'v':
        return f(a[1], d)
    return [a[0]] + [walk(b, f, d + ({'i': [0, 0, 1], 'q': [1], 'd': [2]}.get(a[0], [0] * len(a)))[j])
                     for j, b in enumerate(a[1:])]


def shift(a, k=1):
    return walk(a, lambda i, d: ['v', i + k if i >= d else i])


def sub(a, t, drop=1):
    return walk(a, lambda i, d: ['v', i] if i < d else shift(t, d) if i == d else ['v', i - drop])


def infer(p, g, n):
    match p:
        case ['hyp', i]:
            require(type(i) is int and 0 <= i < len(g))
            return g[i]
        case ['zero', a, h]:
            term(a, n)
            term(h, n + 1)
            return ['=', ['i', ['z'], a, h], a]
        case ['succ', c, a, h]:
            term(c, n)
            term(a, n)
            term(h, n + 1)
            return ['=', ['i', ['s', c], a, h], sub(h, ['i', c, a, h])]
        case ['qe', t, p, q]:
            term(t, n)
            a = infer(p, g, n)
            require(a[0] == 'q')
            require(infer(q, g, n) == sub(a[1], ['z']))
            return sub(a[1], t)
        case ['de', a, p, q]:
            formula(a, n)
            b = infer(p, g, n)
            require(b[0] == 'd')
            require(infer(q, [['q', b[1]]] + [shift(x) for x in g], n + 1) == shift(a))
            return a
        case ['di', b, t, p]:
            formula(b, n + 2)
            term(t, n)
            a = sub(['q', b], t)[1]
            h = a if a[0] == '=' and a[1][:3] == a[2][:3] == ['i', ['v', 0], ['v', 0]] else shift(sub(a, ['z']))
            require(infer(p, [h] + [shift(x) for x in g], n + 1) == sub(a, ['s', ['v', 0]], 0))
            return ['d', b]
        case ['tr', a, b, p, q]:
            term(a, n + 1)
            term(b, n + 1)
            e = infer(p, g, n)
            require(e[0] == '=')
            s, t = e[1:]
            require(infer(q, g, n) == ['=', sub(a, s), sub(b, s)])
            return ['=', sub(a, t), sub(b, t)]
        case _:
            raise ValueError('Invalid proof')


def check(c):
    require(type(c) is dict and set(c) == {'variables', 'context', 'goal', 'proof'})
    n, g, a, p = (c[k] for k in ['variables', 'context', 'goal', 'proof'])
    require(type(n) is int and n >= 0 and type(g) is list)
    for b in g:
        formula(b, n)
    formula(a, n)
    require(infer(p, g, n) == a)


if __name__ == '__main__':
    try:
        check(json.load(sys.stdin))
    except (ValueError, TypeError, IndexError, KeyError, RecursionError):
        sys.exit(1)
    print('ACCEPT')
\end{lstlisting}

\begin{thebibliography}{10}
\bibitem{repo} ProveIt contributors, \emph{Concise presentations of
intuitionistic primitive recursive arithmetic}, README at commit
\texttt{6bf83d8765614b0af51b35747e07fb057d6b2f42}.
\url{https://github.com/VladimirReshetnikov/ProveIt/tree/6bf83d8765614b0af51b35747e07fb057d6b2f42/Logic/PrimitiveRecursiveArithmetic/ConciseIPRA}.
\bibitem{native} ProveIt contributors, \emph{Native.lean} and
\emph{NativeStruct.lean}, same commit, \texttt{Lean/ConciseIPRA/}.
\bibitem{bridge} ProveIt contributors, \emph{NativeBiInterpretation.lean}
and \emph{BiInterpretation.lean}, same commit, \texttt{Lean/ConciseIPRA/}.
\bibitem{witness} ProveIt contributors, \emph{witness\_rule} research
package and \emph{NativeOneIntro.lean}, same commit.
\bibitem{guarded} ProveIt contributors, \emph{guarded\_core} research
package, same commit. See also \emph{guarded\_source} and
\emph{deterministic\_auxiliary}.
\bibitem{uniform} ProveIt contributors, \emph{uniform\_progress} research
package, same commit.
\bibitem{erasure} ProveIt contributors, inherited
\emph{scalar\_predecessor.tex},
\path{Research/certified_zero/reference/}, same commit.
\bibitem{forward} ProveIt contributors, inherited
\emph{equality\_reflection.tex},
\path{Research/certified_zero/reference/}, same commit.
\bibitem{progress} ProveIt contributors, inherited
\emph{primitive\_recursion.tex},
\path{Research/certified_zero/reference/}, same commit.
\bibitem{salehi} S. Salehi, \emph{Decidable formulas of intuitionistic
primitive recursive arithmetic}, Reports on Mathematical Logic 36
(2002), 55--61.
\url{https://rml.tcs.uj.edu.pl/rml-36/36-SALEHI.pdf}.
\end{thebibliography}
\end{document}
